\PassOptionsToPackage{table}{xcolor}
\documentclass[letterpaper,openany]{book}

\input{Preamble.tex}
% \usepackage[utf8]{inputenc}

\begin{document}
\include{TitlePage.tex}

\begingroup
\renewcommand{\baselinestretch}{1.0}
% \tableofcontents
% \listoffigures
% \listoftables
\endgroup

\chapter{Introduction}
{\em dwmwhat(1)} is a command line utility that is very similar to {\em what(1)}
from SCCS.  It searches one or more files for strings that start with "@(\#)"
and displays them on standard output.  Such strings can be embedded in
binaries to represent version information.

The {\em dwmwhat} package includes a C++ class template \texttt{Dwm::What::Info}
that can be used to create structured versions of these strings from string
literals.  This class template is intended to be instantiated to identify
the package name, version and other information inside a string that will be
visible in compiled code.  It is constructed at compile time, and because it
is structured, {\em dwmwhat(1)} can parse it and hence present it decomposed
into its components (as it does when the \texttt{-j} command line option is
used).

For example, let's say you create a trivial program in \texttt{myprog.cc}
as seen in \autoref{lst:myprog}:

\begin{small}
\begin{lstlisting}[language=c++,caption={\texttt{myprog.cc}},captionpos=b,label={lst:myprog}]
#include <iostream>
#include "DwmWhatInfo.hh"

//----------------------------------------------------------------------------
Dwm::What::Info  myproginfo(DWM_WHAT_TYPE_EXE, DWM_WHAT_STATUS_REL, "myprog",
                            "1.0.0", DWM_WHAT_SYM_COPYRIGHT " Jane Doe 2025",
                            "hello!");

//----------------------------------------------------------------------------
int main(int argc, char *argv[])
{
  std::cout << myproginfo.data_view() << '\n';
  std::cout << myproginfo.as_json() << '\n';
  return 0;
}
\end{lstlisting}
\end{small}

If you compile and link this program, {\em dwmwhat(1)} will report this:

\begin{ttfamily}
\% dwmwhat myprog \\
\emoji{robot} \emoji{white-check-mark} myprog 1.0.0 \copyright\ Jane Doe 2025 hello! \\
\emoji{robot} \# \emoji{white-check-mark} dwmwhat 1.0.0 \copyright\ Daniel McRobb 2025 mcplex.net
\end{ttfamily}

And can also report in JSON form:

\begin{ttfamily}
\% dwmwhat -j myprog\\
\parbox{\linewidth}{[\{"file":"myprog","infos":[\{"copyright":"\copyright\ Jane Doe 2025","name":"myprog", "other":"hello!","status":"\emoji{white-check-mark}","type":"\emoji{robot}","version":"1.0.0", "view":"@(\#) \emoji{robot} \emoji{white-check-mark} myprog 1.0.0 \copyright\ Jane Doe 2025 hello!"\},\{"copyright": "\copyright\ Daniel McRobb 2025","name":"dwmwhat","other":"mcplex.net","status":"\emoji{white-check-mark}", "type":"\emoji{robot}\#","version":"1.0.0","view":"@(\#) \emoji{robot}\# \emoji{white-check-mark} dwmwhat 1.0.0 \copyright\ Daniel McRobb 2025 mcplex.net"\}]\}]}
\end{ttfamily}

\section{Under the hood}
Let's look at what's going on under the hood.

We're building a null-terminated string literal at compile time (the
\texttt{Dwm::What::Info} constructor is \texttt{consteval}).  Note that
the only reason we need the null termination is to allow {\em dwmwhat(1)}
and similar tools to find the end of the string in a binary.

The constructor call:

\begin{verbatim}
   Dwm::What::Info(DWM_WHAT_TYPE_EXE, DWM_WHAT_STATUS_REL, "myprog", "1.0.0",
                   DWM_WHAT_SYM_COPYRIGHT " Jane Doe 2025", "hello!");
\end{verbatim}

Produces a string as shown in \autoref{fig:myprogString01}.

\begin{minipage}{\linewidth}
\setcounter{lastbyte}{1}

\begin{center}
  \noindent\begin{tikzpicture}[
    start chain = A going right,
    node distance = 0pt, % Ensure no gap between nodes
    byte node/.style = {draw=lightgray, minimum height=20pt, minimum width=(\textwidth)/56, text width=\dimexpr(\textwidth)/56\relax, text height=16pt, text depth=4pt, font=\ttfamily, anchor=south, align=center, inner sep=0pt, outer sep=0pt, on chain=A}
]
  \def\mypfxstring{{@},{(},{\#},{)}}
  \def\mydelimstring{0xC2,0xA0}
  \def\mysegonestring{0xF0,0x9F,0xA4,0x96}
  \def\mysegtwostring{0xE2,0x9C,0x85}
  \def\mysegthreestring{m,y,p,r,o,g}
  \def\mysegfourstring{1,.,0,.,0}
  \def\copyrightstring{0xC2,0xA9}
  \def\mysegfivestring{{ },J,a,n,e,{ },D,o,e,{ },2,0,2,5}
  \def\mysegsixstring{h,e,l,l,o,!}

  \foreach \byte [count=\i] in \mypfxstring {
    \node[byte node, fill=white] (byte-\arabic{lastbyte}) {\byte};
    \stepcounter{lastbyte}
  }
  \draw [decorate, decoration={brace, amplitude=10pt}]
    (byte-1.north west) -- (byte-4.north east)
    node [above=8pt, midway] {@(\#)};
  \foreach \byte [count=\i] in \mydelimstring {
    \node[byte node, fill=red!15, font=\scriptsize\ttfamily, text depth=0pt] (byte-\arabic{lastbyte}) {\rotatebox{90}{\byte}};
    \stepcounter{lastbyte}
  }
  \draw [decorate, decoration={brace, amplitude=10pt}]
  (byte-5.north west) -- (byte-6.north east)
  node [above=8pt, midway] {\textvisiblespace};

  \foreach \byte [count=\i] in \mysegonestring {
    \node[byte node, fill=white, font=\scriptsize\ttfamily, text depth=0pt] (byte-\arabic{lastbyte}) {\rotatebox{90}{\byte}};
    \stepcounter{lastbyte}
  }
  \draw [decorate, decoration={brace, amplitude=10pt}]
    (byte-7.north west) -- (byte-10.north east)
    node [above=8pt, midway] {\emoji{robot}};
  \foreach \byte [count=\i] in \mydelimstring {
    \node[byte node, fill=red!15, font=\scriptsize\ttfamily, text depth=0pt] (byte-\arabic{lastbyte}) {\rotatebox{90}{\byte}};
    \stepcounter{lastbyte}
  }
  \draw [decorate, decoration={brace, amplitude=10pt}]
  (byte-11.north west) -- (byte-12.north east)
  node [above=8pt, midway] {\textvisiblespace};

  \foreach \byte [count=\i] in \mysegtwostring {
    \node[byte node, fill=white, font=\scriptsize\ttfamily, text depth=0pt] (byte-\arabic{lastbyte}) {\rotatebox{90}{\byte}};
    \stepcounter{lastbyte}
  }
  \draw [decorate, decoration={brace, amplitude=10pt}]
  (byte-13.north west) -- (byte-15.north east)
  node [above=8pt, midway] {\emoji{white-check-mark}};
  \foreach \byte [count=\i] in \mydelimstring {
    \node[byte node, fill=red!15, font=\scriptsize\ttfamily, text depth=0pt] (byte-\arabic{lastbyte}) {\rotatebox{90}{\byte}};
    \stepcounter{lastbyte}    
  }
  \draw [decorate, decoration={brace, amplitude=10pt}]
  (byte-16.north west) -- (byte-17.north east)
  node [above=8pt, midway] {\textvisiblespace};

  \foreach \byte [count=\i] in \mysegthreestring {
    \node[byte node, fill=white] (byte-\arabic{lastbyte}) {\byte};
    \stepcounter{lastbyte}
  }
  \draw [decorate, decoration={brace, amplitude=10pt}]
    (byte-18.north west) -- (byte-23.north east)
    node [above=8pt, midway] {myprog};
  \foreach \byte [count=\i] in \mydelimstring {
    \node[byte node, fill=red!15, font=\scriptsize\ttfamily, text depth=0pt] (byte-\arabic{lastbyte}) {\rotatebox{90}{\byte}};
    \stepcounter{lastbyte}    
  }
  \draw [decorate, decoration={brace, amplitude=10pt}]
  (byte-24.north west) -- (byte-25.north east)
  node [above=8pt, midway] {\textvisiblespace};

  \foreach \byte [count=\i] in \mysegfourstring {
    \node[byte node, fill=white] (byte-\arabic{lastbyte}) {\byte};
    \stepcounter{lastbyte}
  }
  \draw [decorate, decoration={brace, amplitude=10pt}]
    (byte-26.north west) -- (byte-30.north east)
    node [above=8pt, midway] {1.0.0};
  \foreach \byte [count=\i] in \mydelimstring {
    \node[byte node, fill=red!15, font=\scriptsize\ttfamily, text depth=0pt] (byte-\arabic{lastbyte}) {\rotatebox{90}{\byte}};
    \stepcounter{lastbyte}    
  }
  \draw [decorate, decoration={brace, amplitude=10pt}]
  (byte-31.north west) -- (byte-32.north east)
  node [above=8pt, midway] {\textvisiblespace};

  \foreach \byte [count=\i] in \copyrightstring {
    \node[byte node, fill=white, font=\scriptsize\ttfamily, text depth=0pt] (byte-\arabic{lastbyte}) {\rotatebox{90}{\byte}};
    \stepcounter{lastbyte}    
  }
  \foreach \byte [count=\i] in \mysegfivestring {
    \node[byte node, fill=white] (byte-\arabic{lastbyte}) {\byte};
    \stepcounter{lastbyte}
  }
 \draw [decorate, decoration={brace, amplitude=10pt}]
   (byte-33.north west) -- (byte-48.north east)
   node [above=8pt, midway] {\copyright\ Jane Doe 2025};
  \foreach \byte [count=\i] in \mydelimstring {
    \node[byte node, fill=red!15, font=\scriptsize\ttfamily, text depth=0pt] (byte-\arabic{lastbyte}) {\rotatebox{90}{\byte}};
    \stepcounter{lastbyte}    
  }
 \draw [decorate, decoration={brace, amplitude=10pt}]
   (byte-49.north west) -- (byte-50.north east)
   node [above=8pt, midway] {\textvisiblespace};

  \foreach \byte [count=\i] in \mysegsixstring {
    \node[byte node, fill=white] (byte-\arabic{lastbyte}) {\byte};
    \stepcounter{lastbyte}
  }
  \draw [decorate, decoration={brace, amplitude=10pt}]
    (byte-51.north west) -- (byte-56.north east)
    node [above=8pt, midway] {hello!};

  \node[byte node, fill=lightgray!50] (termnull) {$\emptyset$};
  
\end{tikzpicture}
\captionof{figure}{\texttt{Dwm::What::Info} layout}
\label{fig:myprogString01}
\end{center}
\end{minipage}


All we've done here is concatenate {\em segments} of string literals,
separated by a fixed {\em delimiter}.  Some of our string literals are
Unicode code points in UTF-8, such as \emoji{robot}
(\texttt{"\textbackslash{xF0}\textbackslash{x9F}\textbackslash{xA4}\textbackslash{x96}"})
and \emoji{white-check-mark}
(\texttt{"\textbackslash{xE2}\textbackslash{x9C}\textbackslash{x85}"}),
for which I created macros in \texttt{DwmWhatInfo.hh}.  And for the
copyright segment, I concatenate \texttt{DWM\_WHAT\_SYM\_COPYRIGHT}
(\texttt{"\textbackslash{xC2}\textbackslash{xA9}"}) with " Jane Doe 2025".

\subsection{Choice of delimiter}
The delimiter is key to being able to decompose the string in tools
like {\em dwmwhat(1)}.  It needs to be distinct, and never appear inside
a {\em segment}.

It turns out that this was probably the most time-consuming part of
doing this work.  Let's start with what I wanted in a delimiter.

\subsubsection{Desirable delimiter features}
\begin{enumerate}
\item Painless to disallow in the user-specified content ({\em segments}).
We want something that the user will never need.
\item Does not clutter the display of the entire literal
(including the delimiters).
\item Human readability of strings ({\em segments} concatenated with
{\em delimiters}, together) in binary files.
\item Machine readability of strings in binary files (the
ability to decompose into segments), without compromising the
human readability.
\end{enumerate}

\section{Delimiter is Unicode {\em No-Break Space (NBSP)} (U+00A0)}
I’ve chosen Unicode {\em No-Break Space (NBSP)} as the delimiter,
since it serves the needs well. For a human it just looks like a
single space, while it can be distinguished from an ASCII space
(\texttt{0x20}) when parsing.  The UTF-8 encoding for {\em NBSP} is
\texttt{"\textbackslash{xC2}\textbackslash{xA0}"}, i.e.
\texttt{\{0xC2,0xA0\}}.

When you want single spaces in any of your segments, just use ASCII
spaces (\texttt{0x20}) and not {\em NBSP}.

\chapter{Typical Use}
\section{Executable}
In an executable, just declare a static constexpr
\texttt{Dwm::What::Info} in the same translation unit as
\texttt{main()}, preferably in a unique namespace.  You may want to
mark your instance with \texttt{\_\_attribute\_\_((used))} to prevent
the toolchain from optimizing it out of existence due to lack of use.
For example:

\begin{verbatim}
namespace MyNamespace {
   namespace myproginfo {
      static constexpr __attribute((used))
          info(DWM_WHAT_TYPE_EXE, DWM_WHAT_STATUS_REL, "myprog", "1.0.0",
               @DWM_WHAT_SYM_COPYRIGHT "My Name 2025", "my program");
   }
}
\end{verbatim}

\section{Compiled library}
In one of your source files, just declare a static constexpr
\texttt{Dwm::What::Info}, preferably in a unique namespace.  You may
want to mark your instance with \texttt{\_\_attribute\_\_((used))}
to prevent the toolchain from optimizing it out of existence due to
lack of use.  For example:

\begin{verbatim}
namespace MyNamespace {
   namespace mylibinfo {
      static constexpr  __attribute__((used))
         info(DWM_WHAT_TYPE_LIB, DWM_WHAT_STATUS_REL, "mylib", "1.0.0",
              @DWM_WHAT_SYM_COPYRIGHT "My Name 2025", "my library");
   }
}
\end{verbatim}

\section{Header-only library}
In one of your header files, declare an inline constexpr
\texttt{Dwm::What::Info}, preferably in a unique namespace.  You may
want to mark your instance with \texttt{\_\_attribute\_\_((used))}
to prevent the toolchain from optimizing it out of existence due to
lack of use.  For example:

\begin{verbatim}
namespace MyNamespace {
   namespace myhdrinfo {
      inline constexpr  __attribute__((used))
         info(DWM_WHAT_TYPE_HDR, DWM_WHAT_STATUS_REL, "myheaderlib", "1.0.0",
              @DWM_WHAT_SYM_COPYRIGHT "My Name 2025", "my header library");
   }
}
\end{verbatim}

\subsection{Caveats}
It's worth noting that the usual caveats apply here.  Despite the fact that
this is the recommended means of producing a single instance in a linked
program, ODR (One Definition Rule) violations can rear their head in
difficult-to-diagnose ways.  Your header is making a promise it can't
really keep (unless the header never changes), since each compilation
unit is going to get its own copy and the linker is going to pick one.
If you happen to have an object file in a old library that was compiled
against an older version, and link it to an object file that was compiled
against a newer version, the linker is going to just pick one, with
NDR (No Diagnostic Required) for ODR violations.  My recommendation is
to name your instance differently for each release, and provide an
accessor that returns a reference to that instance.  Then the linker
will see two identifiers instead of one, and will keep both.

I mention this because in the case of
\texttt{Dwm::What::SegmentedLiteral} and \texttt{Dwm::What::Info}, a
declaration is typically not just stamping out a variable, but a new
type.  The template is parameterized by the number and size of the
segments, so any change to either of these will yield a different
type, not just different content.  The linkers of today don't seem to
use this information to flag ODR violations.


\chapter{Class templates in \texttt{Dwm::What} namespace}
There are two class templates of interest in the {\em dwmwhat} package:
\texttt{Dwm::What::SegmentedLiteral} and \texttt{Dwm::What::Info}.
We saw simple use of the \texttt{Dwm::What::Info} class template in
the introduction, and it is likely the only class template most users
will use.

However, \texttt{Dwm::What::Info} inherits from
\texttt{Dwm::What::SegmentedLiteral}, and I'll introduce them in
base to derived order.  
\texttt{Dwm::What::SegmentedLiteral} has no constraints on number
of segments or choice of delimiter (an empty delimiter is permissible),
and all of the segments are user defined.

\section{\texttt{Dwm::What::SegmentedLiteral}}
\texttt{Dwm::What::SegmentedLiteral} builds a string literal from
a given delimiter (placed between segments) and one or more segments
(given as string literals).  All arguments to the constructor are
string literals, and used to deduce the template parameters.

An empty delimiter is allowed, as are empty segments.

\subsection{Synopsis}

\begin{verbatim}
template <std::size_t DelimLen, std::size_t FirstLen, std::size_t ...SegLen>
class SegmentedLiteral
{
public:
   consteval SegmentedLiteral(const char (&delim)[DelimLen],
                              const char (&firstSeg)[FirstLen],
                              const char (&...moreSegs)[SegLen]);
   constexpr std::string_view view() const noexcept;
   constexpr std::size_t num_segments() const noexcept;      
   constexpr std::string_view nth(std::size_t n) const noexcept;
};
\end{verbatim}

\subsection{Member functions}

\begin{minipage}{\linewidth}
\subsubsection{Constructor}
The \texttt{Dwm::What::SegmentedLiteral} constructor requires a
delimiter (a string literal) and at least one segment.  All segments
are given as string literals.

\begin{description}[noitemsep,align=left,leftmargin=*,labelsep=10pt,itemindent=-\leftmargin,style=nextline]
  \item[\texttt{delim}] The delimiter that will be placed between segments,
  which may be empty.
  \item[\texttt{firstSeg}] The first segment.
  \item[\texttt{moreSegs}] The remaining segments.
\end{description}
\end{minipage}

\subsubsection{Accessors}
\begin{description}[noitemsep,align=left,leftmargin=*,labelsep=10pt,itemindent=-\leftmargin,style=nextline]
  \item[\texttt{view()}] returns a \texttt{std::string\_view} of the entire
  constructed string literal, sans the terminating null.
  \item[\texttt{num\_segments()}] returns the number of segments.
  \item[\texttt{nth(std::size\_t n)}] returns the \texttt{n}th segment (0 indexed).
\end{description}

\subsection{Examples}

\begin{minipage}{\linewidth}
\input{SegmentedLiteralEx01.tex}
\end{minipage}

\begin{verbatim}
Dwm::What::SegmentedLiteral("|", "@(#)", "DwmWhat", "1.0.0",
                            "Copyright Daniel McRobb 2025", "mcplex.net");
\end{verbatim}

\begin{center}
\noindent\begin{tikzpicture}[
    start chain = A going right,
    node distance = 0pt, % Ensure no gap between bit nodes
    byte node/.style = {draw=lightgray, minimum height=20pt, minimum width=(\textwidth)/56, text width=\dimexpr(\textwidth)/56\relax, text height=16pt, text depth=4pt, font=\ttfamily, anchor=south, align=center, inner sep=0pt, outer sep=0pt, on chain=A}
  ]
  \def\mybytestring{{@},{(},{\#},{)},{|},D,w,m,W,h,a,t,{|},1,.,0,.,0,{|},C,o,p,y,r,i,g,h,t,{ },D,a,n,i,e,l,{ },M,c,R,o,b,b,{ },2,0,2,5,{|},m,c,p,l,e,x,.,n,e,t}
  % Draw the byte nodes and indices
  \setcounter{lastbyte}{0}
  \foreach \byte [count=\i] in \mybytestring {
    \node[byte node, fill=white] (byte-\i) {\byte};
    \stepcounter{lastbyte}
  }
  \node[byte node, fill=lightgray!50] (termnull) {$\emptyset$};
  
  % and a label for the byte nodes
  \ifnum \value{lastbyte} > 0
    \draw [decorate, decoration={brace, amplitude=10pt}] (byte-1.north west) -- (byte-\thelastbyte.north east)
    node [above=8pt, midway, font=\scriptsize] {\texttt{view()}};
    \draw [decorate, decoration={brace, mirror, amplitude=10pt}] (byte-1.south west) -- (byte-4.south east)
    node [below=8pt, midway, font=\scriptsize] {\texttt{nth(0)}};
    \draw [decorate, decoration={brace, mirror, amplitude=10pt}] (byte-6.south west) -- (byte-12.south east)
    node [below=8pt, midway, font=\scriptsize] {\texttt{nth(1)}};
    \draw [decorate, decoration={brace, mirror, amplitude=10pt}] (byte-14.south west) -- (byte-18.south east)
    node [below=8pt, midway, font=\scriptsize] {\texttt{nth(2)}};
    \draw [decorate, decoration={brace, mirror, amplitude=10pt}] (byte-20.south west) -- (byte-47.south east)
    node [below=8pt, midway, font=\scriptsize] {\texttt{nth(3)}};
    \draw [decorate, decoration={brace, mirror, amplitude=10pt}] (byte-49.south west) -- (byte-58.south east)
    node [below=8pt, midway, font=\scriptsize] {\texttt{nth(4)}};

  \draw [decorate, decoration={brace, mirror, amplitude=10pt}] (byte-5.south west) -- (byte-5.south east) 
  node [below=8pt, midway, rotate=270, anchor=east, xshift=3em, yshift=.5pt, font=\scriptsize] {delim};
  \draw [decorate, decoration={brace, mirror, amplitude=10pt}] (byte-13.south west) -- (byte-13.south east) 
  node [below=8pt, midway, rotate=270, anchor=east, xshift=3em, yshift=.5pt, font=\scriptsize] {delim};
  \draw [decorate, decoration={brace, mirror, amplitude=10pt}] (byte-19.south west) -- (byte-19.south east) 
  node [below=8pt, midway, rotate=270, anchor=east, xshift=3em, yshift=.5pt, font=\scriptsize] {delim};
  \draw [decorate, decoration={brace, mirror, amplitude=10pt}] (byte-48.south west) -- (byte-48.south east) 
  node [below=8pt, midway, rotate=270, anchor=east, xshift=3em, yshift=.5pt, font=\scriptsize] {delim};

  \fi
\end{tikzpicture}
\captionof{figure}{\texttt{Dwm::What::SegmentedLiteral} layout}
\end{center}


\pagebreak

\section{\texttt{Dwm::What::Info}}
\texttt{Dwm::What::Info} is just a \texttt{Dwm::What::SegmentedLiteral}
with seven segments and a specific delimiter (not user-specified).
The first segment (\texttt{"@(\#)"}) is automatically added, and named
accessors are provided for the remaining six segments.

\begin{minipage}{\linewidth}
\subsection{Synopsis}

\begin{small}
\begin{verbatim}
template <std::size_t TypeLen, std::size_t StatusLen, std::size_t NameLen,
           std::size_t VersionLen, std::size_t CopyrightLen, std::size_t OtherLen>
class Info
   : public SegmentedLiteral<sizeof(DWM_WHAT_DELIM), sizeof("@(#)"), TypeLen,
                             StatusLen, NameLen, VersionLen, CopyrightLen, OtherLen>
{
public:
   consteval Info(const char (&type)[TypeLen], const char (&status)[StatusLen],
                  const char (&name)[NameLen], const char (&version)[VersionLen],
                  const char (&copyright)[CopyrightLen], const char (&other)[OtherLen]);

   constexpr std::string_view type() const noexcept;
   constexpr std::string_view status() const noexcept;
   constexpr std::string_view name() const noexcept;
   constexpr std::string_view version() const noexcept;
   constexpr std::string_view copyright() const noexcept;
   constexpr std::string_view other() const noexcept;
   constexpr std::string_view data_view() const noexcept;
   constexpr std::string_view view() const noexcept;
};
\end{verbatim}
\end{small}
\end{minipage}

\subsection{Constructor}
\subsubsection{constructor arguments}
\begin{description}[noitemsep,align=left,leftmargin=*,labelsep=10pt,itemindent=-\leftmargin,style=nextline]
\item[\texttt{type}] The type of the package.  By convention, one or more of
  \texttt{DWM\_WHAT\_TYPE\_EXE}, \texttt{DWM\_WHAT\_TYPE\_LIB},
  \texttt{DWM\_WHAT\_TYPE\_HDR} and \texttt{DWM\_WHAT\_TYPE\_DOC}.  These are
  just macros that expand to string literals.  See \autoref{table:Macros}.
\item[\texttt{status}] The status of the package.  By convention, one of
  \texttt{DWM\_WHAT\_STATUS\_REL}, \texttt{DWM\_WHAT\_STATUS\_RC} or
  \texttt{DWM\_WHAT\_STATUS\_DEV}.  These are just macros that expand to
  string literals.  See \autoref{table:Macros}.
  \item[\texttt{name}] The name of the package.
  \item[\texttt{version}] The version of the package.
  \item[\texttt{copyright}] The copyright.
  \item[\texttt{other}] Any other additional information.
\end{description}

\begin{minipage}{\linewidth}
\subsection{Accesssors}
\begin{description}[noitemsep,align=left,leftmargin=*,labelsep=10pt,itemindent=-\leftmargin,style=nextline]
\item[\texttt{type()}] Returns the type of the package.  This is the same that was
  provided as the first argument to the constructor.
\item[\texttt{status()}] Returns the status of the package.  This is the same that was
  provided as the second argument to the constructor.
\item[\texttt{name()}] Returns the name of the package.  This is the same that was
  provided as the third argument to the constructor.
\item[\texttt{version()}] Returns the version of the package.  This is the same that
  was provided as the fourth argument to the constructor.
\item[\texttt{copyright()}] Returns the copyright.  This is the same that was provided
  as the fifth argument to the constructor.
\item[\texttt{other()}] Returns any additional information that was provided.  This
  is the same as the sixth (final) argument to the constructor.
\end{description}
\end{minipage}

\subsection{Equivalence to \texttt{Dwm::What::SegmentedLiteral}}
\begin{minipage}{\linewidth}
The following will produce the equivalent string literal.  This makes
it obvious that \texttt{Dwm::What::Info} is just a
\texttt{Dwm::What::SegmentedLiteral} with a specific delimiter
(\texttt{"\textbackslash{xC2}\textbackslash{xA0}"}) and an
automatically added first segment (\texttt{"@(\#)"}).

\begin{Verbatim}
   SegmentedLiteral(\textcolor{blue}{"\textbackslash{xC2}\textbackslash{xA0}"}, \textcolor{blue}{"@(\#)"},
                    DWM_WHAT_TYPE_LIB, DWM_WHAT_STATUS_REL, "DwmWhat",
                    "1.0.0", "Daniel McRobb 2025", "mcplex.net");

               Info(DWM_WHAT_TYPE_LIB, DWM_WHAT_STATUS_REL, "DwmWhat",
                    "1.0.0", "Daniel McRobb 2025", "mcplex.net");
\end{Verbatim}
\end{minipage}

\subsection{Examples}
\begin{minipage}{\linewidth}
This call to the \texttt{Dwm::What::Info} constructor:
\begin{verbatim}
   Dwm::What::Info(DWM_WHAT_TYPE_EXE DWM_WHAT_TYPE_HDR, DWM_WHAT_STATUS_REL,
                   "DwmWhat", "1.0.0", "Daniel McRobb 2025", "mcplex.net");
\end{verbatim}

Would produce a 67 byte string literal at compile time with the layout shown
in \autoref{figure:InfoEx01}.  Note this figure shows what would be returned
by accessors inherited from the \texttt{Dwm::What::SegmentedLiteral} base
class (in black) and the accessors from the \texttt{Dwm::What::Info} class
(in \textcolor{blue}{blue}).

\begin{center}
  \begin{figure}[H]
\noindent\begin{tikzpicture}[
  start chain = A going right,
  node distance = 0pt, % Ensure no gap between bit nodes
  byte node/.style = {draw=lightgray, minimum height=20pt, text height=16pt, text depth=4pt, font=\ttfamily, anchor=south, align=center, inner sep=0pt, outer sep=0pt, on chain=A}
]
\newlength{\onebyte}
\setlength{\onebyte}{1.7ex}

\node[byte node,minimum width=4\onebyte,text width=4\onebyte,fill=white] (seg1) {@(\#)};
\node[byte node,minimum width=\onebyte,text width=\onebyte, fill=red!15,font=\scriptsize\ttfamily, text depth=0pt] (delim1a) {\rotatebox{90}{0xC2}};
\node[byte node,minimum width=\onebyte,text width=\onebyte, fill=red!15,font=\scriptsize\ttfamily, text depth=0pt] (delim1b) {\rotatebox{90}{0xA0} };
\node[byte node,minimum width=7\onebyte,text width=7\onebyte,fill=white] (seg2) {\emoji{robot}\#};
\node[byte node,minimum width=\onebyte,text width=\onebyte, fill=red!15,font=\scriptsize\ttfamily, text depth=0pt] (delim2a) {\rotatebox{90}{0xC2}};
\node[byte node,minimum width=\onebyte,text width=\onebyte, fill=red!15,font=\scriptsize\ttfamily, text depth=0pt] (delim2b) {\rotatebox{90}{0xA0} };
\node[byte node,minimum width=3\onebyte,text width=3\onebyte,fill=white] (seg3) {\emoji{white-check-mark}};
\node[byte node,minimum width=\onebyte,text width=\onebyte, fill=red!15,font=\scriptsize\ttfamily, text depth=0pt] (delim3a) {\rotatebox{90}{0xC2}};
\node[byte node,minimum width=\onebyte,text width=\onebyte, fill=red!15,font=\scriptsize\ttfamily, text depth=0pt] (delim3b) {\rotatebox{90}{0xA0} };
\node[byte node,minimum width=7\onebyte,text width=7\onebyte,fill=white] (seg4) {DwmWhat};
\node[byte node,minimum width=\onebyte,text width=\onebyte, fill=red!15,font=\scriptsize\ttfamily, text depth=0pt] (delim4a) {\rotatebox{90}{0xC2}};
\node[byte node,minimum width=\onebyte,text width=\onebyte, fill=red!15,font=\scriptsize\ttfamily, text depth=0pt] (delim4b) {\rotatebox{90}{0xA0} };
\node[byte node,minimum width=5\onebyte,text width=5\onebyte,fill=white] (seg5) {1.0.0};
\node[byte node,minimum width=\onebyte,text width=\onebyte, fill=red!15,font=\scriptsize\ttfamily, text depth=0pt] (delim5a) {\rotatebox{90}{0xC2}};
\node[byte node,minimum width=\onebyte,text width=\onebyte, fill=red!15,font=\scriptsize\ttfamily, text depth=0pt] (delim5b) {\rotatebox{90}{0xA0} };
\node[byte node,minimum width=18\onebyte,text width=18\onebyte,fill=white] (seg6) {Daniel McRobb 2025};
\node[byte node,minimum width=\onebyte,text width=\onebyte, fill=red!15,font=\scriptsize\ttfamily, text depth=0pt] (delim6a) {\rotatebox{90}{0xC2}};
\node[byte node,minimum width=\onebyte,text width=\onebyte, fill=red!15,font=\scriptsize\ttfamily, text depth=0pt] (delim6b) {\rotatebox{90}{0xA0} };
\node[byte node,minimum width=10\onebyte,text width=10\onebyte,fill=white] (seg7) {mcplex.net};
\node[byte node,minimum width=\onebyte,text width=\onebyte,fill=lightgray!50] (termnull) {$\emptyset$};

\node (nth0) [below=of seg1.south, yshift=-20pt,font=\scriptsize]
      {\texttt{nth(0)}};
\draw [decorate, decoration={brace, mirror, amplitude=10pt}]
  (seg2.south west) -- (seg2.south east)
  node [below=8pt, midway, font=\scriptsize, text=blue] {\texttt{type()}};
\node (nth1) [below=of seg2.south, yshift=-20pt,font=\scriptsize]
      {\texttt{nth(1)}};
\draw [decorate, decoration={brace, mirror, amplitude=10pt}]
  (seg3.south west) -- (seg3.south east)
  node [below=8pt, midway, font=\scriptsize, text=blue] {\texttt{status()}};
\node (nth2) [below=of seg3.south, yshift=-20pt,font=\scriptsize]
      {\texttt{nth(2)}};
\draw [decorate, decoration={brace, mirror, amplitude=10pt}]
  (seg4.south west) -- (seg4.south east)
  node [below=8pt, midway, font=\scriptsize, text=blue] {\texttt{name()}};
\node (nth3) [below=of seg4.south, yshift=-20pt,font=\scriptsize]
      {\texttt{nth(3)}};
\draw [decorate, decoration={brace, mirror, amplitude=10pt}]
  (seg5.south west) -- (seg5.south east)
  node [below=8pt, midway, font=\scriptsize, text=blue] {\texttt{version()}};
\node (nth4) [below=of seg5.south, yshift=-20pt,font=\scriptsize]
      {\texttt{nth(4)}};
\draw [decorate, decoration={brace, mirror, amplitude=10pt}]
  (seg6.south west) -- (seg6.south east)
  node [below=8pt, midway, font=\scriptsize, text=blue] {\texttt{copyright()}};
\node (nth6) [below=of seg6.south, yshift=-20pt,font=\scriptsize]
      {\texttt{nth(5)}};
\draw [decorate, decoration={brace, mirror, amplitude=10pt}]
  (seg7.south west) -- (seg7.south east)
  node [below=8pt, midway, font=\scriptsize, text=blue] {other()};
\node (nth8) [below=of seg7.south, yshift=-20pt,font=\scriptsize]
      {\texttt{nth(6)}};
      
\draw [decorate, decoration={brace, amplitude=10pt}]
  (seg2.north west) -- (seg7.north east)
  node [above=8pt, midway, font=\scriptsize, text=blue] {\texttt{data\_view()}};

\draw [decorate, decoration={brace, amplitude=40pt}]
  (seg1.north west) -- (seg7.north east)
  node [above=40pt, midway, font=\scriptsize] {\texttt{view()}};

\end{tikzpicture}
\caption{\texttt{Dwm::What::Info} layout}
\label{figure:InfoEx01}
\end{figure}
\end{center}
\end{minipage}

The 67 bytes:
\begin{tabular}{ll}
  @(\#)            &  4 bytes \\
  delim            &  2 bytes \\
  type             &  7 bytes \\
  delim            &  2 bytes \\
  status           &  3 bytes \\
  delim            &  2 bytes \\
  name             &  7 bytes \\
  delim            &  2 bytes \\
  version          &  5 bytes \\
  delim            &  2 bytes \\
  copyright        & 18 bytes \\
  delim            &  2 bytes \\
  other            & 10 bytes \\
  null             &  1 byte \\ \hline
  \textbf{Total} & \textbf{67 bytes}
\end{tabular}

And {\em dwmwhat} would print this when scanning a compiled program
containing such an instance of \texttt{Dwm::What::Info}:

\texttt{
\emoji{robot}\# \emoji{white-check-mark} DwmWhat 1.0.0 Daniel McRobb 2025 mcplex.net
}


\subsection{\texttt{SegmentedLiteral} and \texttt{Info} equivalents}

\begin{minipage}{\linewidth}
  All of the following will produce the same string, as show in
  \autoref{figure:SLInfoEquiv01}.  The first two are using macros from
  \texttt{DwmWhatInfo.hh} (\autoref{table:Macros}) in constructor
  arguments.  The next two are repeats of the first two, with macro
  substitution.

\begin{small}
\begin{verbatim}
Dwm::What::SegmentedLiteral(DWM_WHAT_SYM_NB_SPACE, "@(#)", DWM_WHAT_TYPE_LIB,
                            DWM_WHAT_STATUS_REL, "dwmwhat", "1.0.0",
                            "Daniel McRobb", "mcplex.net");
\end{verbatim}
\color{blue}
\begin{verbatim}
            Dwm::What::Info(DWM_WHAT_TYPE_LIB, DWM_WHAT_STATUS_REL, "dwmwhat",
                            "1.0.0", "Daniel McRobb", "mcplex.net");
\end{verbatim}
\color{black}

\begin{verbatim}
Dwm::What::SegmentedLiteral("\xC2\xA0", "@(#)", "\xF0\x9F\x93\x9A", "\xE2\x9C\x85",
                            "dwmwhat", "1.0.0", "Daniel McRobb", "mcplex.net");
\end{verbatim}
\color{blue}
\begin{verbatim}
            Dwm::What::Info("\xF0\x9F\x93\x9A", "\xE2\x9C\x85", "dwmwhat",
                            "1.0.0", "Daniel McRobb", "mcplex.net");
\end{verbatim}
\color{black}
\end{small}

\begin{verbatim}
Dwm::What::SegmentedLiteral("\xC2\xA0", "@(#)", DWM_WHAT_TYPE_LIB,
                            DWM_WHAT_STATUS_REL, "DwmWhat", "1.0.0",
                            "Daniel McRobb 2025", "mcplex.net");
\end{verbatim}

\begin{center}
\noindent\begin{tikzpicture}[
    start chain = A going right,
    node distance = 0pt, % Ensure no gap between bit nodes
    byte node/.style = {draw=lightgray, minimum height=20pt, minimum width=(\textwidth)/56, text width=\dimexpr(\textwidth)/56\relax, text height=16pt, text depth=4pt, font=\ttfamily, anchor=south, align=center, inner sep=0pt, outer sep=0pt, on chain=A}
]
  \def\mypfxstring{{@},{(},{\#},{)}}
  \def\mydelimstring{0xC2,0xA0}
  \def\mysegonestring{0xF0,0x9F,0x93,0x9A}
  \def\mysegtwostring{0xE2,0x9C,0x85}
  \def\mysegthreestring{D,w,m,W,h,a,t}
  \def\mysegfourstring{1,.,0,.,0}
  \def\mysegfivestring{D,a,n,i,e,l,{ },M,c,R,o,b,b}
  \def\mysegsixstring{m,c,p,l,e,x,.,n,e,t}

  \foreach \byte [count=\i] in \mypfxstring {
    \node[byte node, fill=white] (byte-\i) {\byte};
    \stepcounter{lastbyte}
  }
  \draw [decorate, decoration={brace, mirror, amplitude=10pt}]
    (byte-1.south west) -- (byte-4.south east)
    node [below=8pt, midway, font=\scriptsize] {\texttt{nth(0)}};
  \draw [decorate, decoration={brace, amplitude=10pt}]
    (byte-1.north west) -- (byte-4.north east)
    node [above=8pt, midway] {@(\#)};
  \foreach \byte [count=\i] in \mydelimstring {
    \node[byte node, fill=red!15, font=\scriptsize\ttfamily, text depth=0pt] (byte-\i) {\rotatebox{90}{\byte}};
  }

  \foreach \byte [count=\i] in \mysegonestring {
    \node[byte node, fill=white, font=\scriptsize\ttfamily, text depth=0pt] (byte-\i) {\rotatebox{90}{\byte}};
    \stepcounter{lastbyte}
  }
  \draw [decorate, decoration={brace, mirror, amplitude=10pt}]
  (byte-7.south west) -- (byte-10.south east)
  node [below=8pt, midway, font=\scriptsize] {\texttt{nth(1)}};
  \draw (byte-7.south west) -- (byte-10.south east)
    node [below=20pt, midway, font=\scriptsize, text=blue] {\texttt{type()}};
  \draw [decorate, decoration={brace, amplitude=10pt}]
  (byte-7.north west) -- (byte-10.north east)
  node [above=8pt, midway] {\emoji{books}};
  \foreach \byte [count=\i] in \mydelimstring {
    \node[byte node, fill=red!15, font=\scriptsize\ttfamily, text depth=0pt] (byte-\i) {\rotatebox{90}{\byte}};
  }

  \foreach \byte [count=\i] in \mysegtwostring {
    \node[byte node, fill=white, font=\scriptsize\ttfamily, text depth=0pt] (byte-\i) {\rotatebox{90}{\byte}};
    \stepcounter{lastbyte}
  }
  \draw [decorate, decoration={brace, mirror, amplitude=10pt}]
  (byte-13.south west) -- (byte-15.south east)
  node [below=8pt, midway, font=\scriptsize] {\texttt{nth(2)}};
  \draw (byte-13.south west) -- (byte-15.south east)
    node [below=20pt, midway, font=\scriptsize, text=blue] {\texttt{status()}};
  \draw [decorate, decoration={brace, amplitude=10pt}]
  (byte-13.north west) -- (byte-15.north east)
  node [above=8pt, midway] {\emoji{white-check-mark}};
  \foreach \byte [count=\i] in \mydelimstring {
    \node[byte node, fill=red!15, font=\scriptsize\ttfamily, text depth=0pt] (byte-\i) {\rotatebox{90}{\byte}};
  }

  \foreach \byte [count=\i] in \mysegthreestring {
    \node[byte node, fill=white] (byte-\i) {\byte};
    \stepcounter{lastbyte}
  }
  \draw [decorate, decoration={brace, mirror, amplitude=10pt}]
    (byte-18.south west) -- (byte-24.south east)
    node [below=8pt, midway, font=\scriptsize] {\texttt{nth(3)}};
  \draw (byte-18.south west) -- (byte-24.south east)
    node [below=20pt, midway, font=\scriptsize, text=blue] {\texttt{name()}};
  \draw [decorate, decoration={brace, amplitude=10pt}]
    (byte-18.north west) -- (byte-24.north east)
    node [above=8pt, midway] {DwmWhat};
  \foreach \byte [count=\i] in \mydelimstring {
    \node[byte node, fill=red!15, font=\scriptsize\ttfamily, text depth=0pt] (byte-\i) {\rotatebox{90}{\byte}};
  }

  \foreach \byte [count=\i] in \mysegfourstring {
    \node[byte node, fill=white] (byte-\i) {\byte};
    \stepcounter{lastbyte}
  }
  \draw [decorate, decoration={brace, mirror, amplitude=10pt}]
    (byte-27.south west) -- (byte-31.south east)
    node [below=8pt, midway, font=\scriptsize] {\texttt{nth(4)}};
  \draw (byte-27.south west) -- (byte-31.south east)
    node [below=20pt, midway, font=\scriptsize, text=blue] {\texttt{version()}};
  \draw [decorate, decoration={brace, amplitude=10pt}]
    (byte-27.north west) -- (byte-31.north east)
    node [above=8pt, midway] {1.0.0};
  \foreach \byte [count=\i] in \mydelimstring {
    \node[byte node, fill=red!15, font=\scriptsize\ttfamily, text depth=0pt] (byte-\i) {\rotatebox{90}{\byte}};
  }

  \foreach \byte [count=\i] in \mysegfivestring {
    \node[byte node, fill=white] (byte-\i) {\byte};
    \stepcounter{lastbyte}
  }
  \draw [decorate, decoration={brace, mirror, amplitude=10pt}]
    (byte-34.south west) -- (byte-46.south east)
    node [below=8pt, midway, font=\scriptsize] {\texttt{nth(5)}};
  \draw (byte-34.south west) -- (byte-46.south east)
    node [below=20pt, midway, font=\scriptsize, text=blue] {\texttt{copyright()}};
  \draw [decorate, decoration={brace, amplitude=10pt}]
    (byte-34.north west) -- (byte-46.north east)
    node [above=8pt, midway] {Daniel McRobb};
  \foreach \byte [count=\i] in \mydelimstring {
    \node[byte node, fill=red!15, font=\scriptsize\ttfamily, text depth=0pt] (byte-\i) {\rotatebox{90}{\byte}};
  }

  \foreach \byte [count=\i] in \mysegsixstring {
    \node[byte node, fill=white] (byte-\i) {\byte};
    \stepcounter{lastbyte}
  }
  \draw [decorate, decoration={brace, mirror, amplitude=10pt}]
    (byte-49.south west) -- (byte-58.south east)
    node [below=8pt, midway, font=\scriptsize] {\texttt{nth(6)}};
  \draw (byte-49.south west) -- (byte-58.south east)
    node [below=20pt, midway, font=\scriptsize, text=blue] {\texttt{other()}};
  \draw [decorate, decoration={brace, amplitude=10pt}]
    (byte-49.north west) -- (byte-58.north east)
    node [above=8pt, midway] {mcplex.net};

  \node[byte node, fill=lightgray!50] (termnull) {$\emptyset$};
    
\end{tikzpicture}
\captionof{figure}{\texttt{Dwm::What::SegmentedLiteral} layout}
\end{center}


\end{minipage}


\chapter{More \texttt{Info} examples}

\section{\texttt{Info} Example 1}
\begin{minipage}{\linewidth}
  
\begin{Verbatim}
Dwm::What::Info(\textcolor{typeColor}{DWM_WHAT_TYPE_LIB}, \textcolor{statusColor}{DWM_WHAT_STATUS_REL}, \textcolor{nameColor}{"mylib"}, \textcolor{versionColor}{"1.0.0"},
                \textcolor{copyrightColor}{"Daniel McRobb"}, \textcolor{otherColor}{"my library"})
\end{Verbatim}
\noindent\begin{tikzpicture}[start chain = A going right]
\node[draw=none,anchor=north west,yshift=1em,on chain=A] {\em compilation produces compile-time string};
\node[draw, single arrow,
  minimum height=10mm, minimum width=15mm,
  single arrow head extend=2mm,
  anchor=east, rotate=-90,xshift=-1em,on chain=A] {};
\end{tikzpicture}

\begin{tikzpicture}[
    start chain = A going right,
    node distance = 0pt, % Ensure no gap between bit nodes
    byte node/.style = {draw=lightgray, minimum height=20pt, minimum width=(\textwidth)/56, text width=\dimexpr(\textwidth)/56\relax, text height=16pt, text depth=4pt, font=\ttfamily, anchor=south, align=center, inner sep=0pt, outer sep=0pt, on chain=A}
]
    
  \def\mypfxstring{{@},{(},{\#},{)}}
  \def\mydelimstring{0xC2,0xA0}
  \def\mysegonestring{0xF0,0x9F,0x93,0x9A}
  \def\mysegtwostring{0xE2,0x9C,0x85}
  \def\mysegthreestring{m,y,l,i,b}
  \def\mysegfourstring{1,.,0,.,0}
  \def\mysegfivestring{D,a,n,i,e,l,{ },M,c,R,o,b,b}
  \def\mysegsixstring{m,y,{ },l,i,b,r,a,r,y}

  \foreach \byte [count=\i] in \mypfxstring {
    \node[byte node, fill=white] (byte-\arabic{lastbyte}) {\byte};
    \stepcounter{lastbyte}
  }
  \draw [decorate, decoration={brace, mirror, amplitude=10pt}]
    (byte-1.south west) -- (byte-4.south east)
    node [below=8pt, midway, font=\scriptsize] {\texttt{nth(0)}};
  \draw [decorate, decoration={brace, amplitude=10pt}]
    (byte-1.north west) -- (byte-4.north east)
    node [above=8pt, midway] {@(\#)};
  \foreach \byte [count=\i] in \mydelimstring {
    \node[byte node, text=delimColor, fill=white, font=\scriptsize\ttfamily, text depth=0pt] (byte-\arabic{lastbyte}) {\rotatebox{90}{\byte}};
    \stepcounter{lastbyte}
  }
  \draw [decorate, decoration={brace, amplitude=10pt}]
  (byte-5.north west) -- (byte-6.north east)
  node [above=8pt, midway] {\textvisiblespace};

  \foreach \byte [count=\i] in \mysegonestring {
    \node[byte node, fill=typeColor!10, text=typeColor, font=\scriptsize\ttfamily, text depth=0pt] (byte-\arabic{lastbyte}) {\rotatebox{90}{\byte}};
    \stepcounter{lastbyte}
  }
  \draw [decorate, decoration={brace, mirror, amplitude=10pt}]
    (byte-7.south west) -- (byte-10.south east)
    node [below=8pt, midway, font=\scriptsize] {\texttt{nth(1)}}
    node [below=18pt, midway, font=\scriptsize, text=typeColor] {\texttt{type()}};
  \draw [decorate, decoration={brace, amplitude=10pt}]
    (byte-7.north west) -- (byte-10.north east)
    node [above=8pt, midway] {\emoji{books}};
  \foreach \byte [count=\i] in \mydelimstring {
    \node[byte node, text=delimColor, fill=white, font=\scriptsize\ttfamily, text depth=0pt] (byte-\arabic{lastbyte}) {\rotatebox{90}{\byte}};
    \stepcounter{lastbyte}
  }
  \draw [decorate, decoration={brace, amplitude=10pt}]
  (byte-11.north west) -- (byte-12.north east)
  node [above=8pt, midway] {\textvisiblespace};

  \foreach \byte [count=\i] in \mysegtwostring {
    \node[byte node, fill=statusColor!10, font=\scriptsize\ttfamily, text=statusColor, text depth=0pt] (byte-\arabic{lastbyte}) {\rotatebox{90}{\byte}};
    \stepcounter{lastbyte}
  }
  \draw [decorate, decoration={brace, mirror, amplitude=10pt}]
    (byte-13.south west) -- (byte-15.south east)
    node [below=8pt, midway, font=\scriptsize] {\texttt{nth(2)}}
    node [below=18pt, midway, font=\scriptsize, text=statusColor] {\texttt{status()}};
  \draw [decorate, decoration={brace, amplitude=10pt}]
  (byte-13.north west) -- (byte-15.north east)
  node [above=8pt, midway] {\emoji{white-check-mark}};
  \foreach \byte [count=\i] in \mydelimstring {
    \node[byte node, text=delimColor, fill=white, font=\scriptsize\ttfamily, text depth=0pt] (byte-\arabic{lastbyte}) {\rotatebox{90}{\byte}};
    \stepcounter{lastbyte}    
  }
  \draw [decorate, decoration={brace, amplitude=10pt}]
  (byte-16.north west) -- (byte-17.north east)
  node [above=8pt, midway] {\textvisiblespace};

  \foreach \byte [count=\i] in \mysegthreestring {
    \node[byte node, fill=nameColor!10, text=nameColor] (byte-\arabic{lastbyte}) {\byte};
    \stepcounter{lastbyte}
  }
  \draw [decorate, decoration={brace, mirror, amplitude=10pt}]
    (byte-18.south west) -- (byte-22.south east)
    node [below=8pt, midway, font=\scriptsize] {\texttt{nth(3)}}
    node [below=18pt, midway, font=\scriptsize, text=nameColor] {\texttt{name()}};
  \draw [decorate, decoration={brace, amplitude=10pt}]
    (byte-18.north west) -- (byte-22.north east)
    node [above=8pt, midway] {mylib};
  \foreach \byte [count=\i] in \mydelimstring {
    \node[byte node, text=delimColor, fill=white, font=\scriptsize\ttfamily, text depth=0pt] (byte-\arabic{lastbyte}) {\rotatebox{90}{\byte}};
    \stepcounter{lastbyte}    
  }
  \draw [decorate, decoration={brace, amplitude=10pt}]
  (byte-23.north west) -- (byte-24.north east)
  node [above=8pt, midway] {\textvisiblespace};

  \foreach \byte [count=\i] in \mysegfourstring {
    \node[byte node, fill=versionColor!10, text=versionColor] (byte-\arabic{lastbyte}) {\byte};
    \stepcounter{lastbyte}
  }
  \draw [decorate, decoration={brace, mirror, amplitude=10pt}]
    (byte-25.south west) -- (byte-29.south east)
    node [below=8pt, midway, font=\scriptsize] {\texttt{nth(4)}}
    node [below=18pt, midway, font=\scriptsize, text=versionColor] {\texttt{version()}};
  \draw [decorate, decoration={brace, amplitude=10pt}]
    (byte-25.north west) -- (byte-29.north east)
    node [above=8pt, midway] {1.0.0};
  \foreach \byte [count=\i] in \mydelimstring {
    \node[byte node, text=delimColor, fill=white, font=\scriptsize\ttfamily, text depth=0pt] (byte-\arabic{lastbyte}) {\rotatebox{90}{\byte}};
    \stepcounter{lastbyte}    
  }
  \draw [decorate, decoration={brace, amplitude=10pt}]
  (byte-30.north west) -- (byte-31.north east)
  node [above=8pt, midway] {\textvisiblespace};

  \foreach \byte [count=\i] in \mysegfivestring {
    \node[byte node, fill=copyrightColor!10, text=copyrightColor] (byte-\arabic{lastbyte}) {\byte};
    \stepcounter{lastbyte}
  }
  \draw [decorate, decoration={brace, mirror, amplitude=10pt}]
    (byte-32.south west) -- (byte-44.south east)
    node [below=8pt, midway, font=\scriptsize] {\texttt{nth(5)}}
    node [below=18pt, midway, font=\scriptsize, text=copyrightColor] {\texttt{copyright()}};
  \draw [decorate, decoration={brace, amplitude=10pt}]
    (byte-32.north west) -- (byte-44.north east)
    node [above=8pt, midway] {Daniel McRobb};
  \foreach \byte [count=\i] in \mydelimstring {
    \node[byte node, text=delimColor, fill=white, font=\scriptsize\ttfamily, text depth=0pt] (byte-\arabic{lastbyte}) {\rotatebox{90}{\byte}};
    \stepcounter{lastbyte}    
  }
  \draw [decorate, decoration={brace, amplitude=10pt}]
  (byte-45.north west) -- (byte-46.north east)
  node [above=8pt, midway] {\textvisiblespace};

  \foreach \byte [count=\i] in \mysegsixstring {
    \node[byte node, fill=otherColor!10, text=otherColor] (byte-\arabic{lastbyte}) {\byte};
    \stepcounter{lastbyte}
  }
  \draw [decorate, decoration={brace, mirror, amplitude=10pt}]
    (byte-47.south west) -- (byte-56.south east)
    node [below=8pt, midway, font=\scriptsize] {\texttt{nth(6)}}
    node [below=18pt, midway, font=\scriptsize, text=otherColor] {\texttt{other()}};
  \draw [decorate, decoration={brace, amplitude=10pt}]
    (byte-47.north west) -- (byte-56.north east)
    node [above=8pt, midway] {my library};

  \node[byte node, fill=lightgray!50] (termnull) {$\emptyset$};

  \draw [decorate, decoration={brace,mirror,amplitude=20pt,raise=24pt}]
    (byte-7.south west) -- (byte-56.south east)
    node [below=40pt, midway, font=\scriptsize,text=blue] {\texttt{data\_view()}};

  \draw [decorate, decoration={brace,mirror,amplitude=20pt,raise=48pt,aspect=0.45}]
    (byte-1.south west) -- (byte-56.south east)
    node [below=64pt, midway, xshift=-2.2em, font=\scriptsize] {\texttt{view()}};
    
\end{tikzpicture}
\label{figure:Info2Seg2StrEx01}

\begin{tikzpicture}[start chain = A going right]
\node[draw=none,minimum height=3em,anchor=north west,yshift=2em,on chain=A] {\em and an array of 7 bytes for segment lengths};
\end{tikzpicture}

\setcounter{lastbyte}{0}
\begin{tikzpicture}[
    start chain = A going below,
    node distance = 0pt, % Ensure no gap between bit nodes
    byte node/.style = {draw=lightgray, minimum height=(\textwidth)/56, minimum width=20pt, text width=16pt, text height=\dimexpr(\textwidth)/56\relax, text depth=2pt, font=\ttfamily, anchor=south, align=center, inner sep=0pt, outer sep=0pt, on chain=A}
]
    
  \def\mysegonestring{0x04,0x04,0x03,0x05,0x05,0x0D,0x0A}

  \foreach \byte [count=\i from 0] in \mysegonestring {
    \node[byte node, font=\scriptsize\ttfamily, text depth=2pt] (byte-\arabic{lastbyte}) {\byte};
    \stepcounter{lastbyte}
  }
  \foreach \byte [count=\i from 0] in \mysegonestring {
    \draw [decorate, decoration={brace, mirror, amplitude=10pt}]
      (byte-\i.north west) -- (byte-\i.south west)
      node [anchor=west, midway, left, xshift=-1em,font=\scriptsize] (seglen-\i) {\texttt{nth(\i).size()}};
  }
  \draw (byte-1.north west) -- (byte-1.south west)
    node [anchor=west,left=of seglen-1,left,font=\scriptsize\ttfamily] {\textcolor{typeColor}{type().size(),}};
  \draw (byte-2.north west) -- (byte-2.south west)
    node [anchor=west,left=of seglen-2,left,font=\scriptsize\ttfamily] {\textcolor{statusColor}{status().size(),}};
  \draw (byte-3.north west) -- (byte-3.south west)
    node [anchor=west,left=of seglen-3,left,font=\scriptsize\ttfamily] {\textcolor{nameColor}{name().size(),}};
  \draw (byte-4.north west) -- (byte-4.south west)
    node [anchor=west,left=of seglen-4,left,font=\scriptsize\ttfamily] {\textcolor{versionColor}{version().size(),}};
  \draw (byte-5.north west) -- (byte-5.south west)
    node [anchor=west,left=of seglen-5,left,font=\scriptsize\ttfamily] {\textcolor{copyrightColor}{copyright().size(),}};
  \draw (byte-6.north west) -- (byte-6.south west)
    node [anchor=west,left=of seglen-6,left,font=\scriptsize\ttfamily] {\textcolor{otherColor}{other().size(),}};

  \draw (byte-0.north west) -- (byte-0.south west)
    node [anchor=east,right=of byte-0,right,font=\scriptsize\ttfamily] {\textcolor{black}{= 4}};
  \draw (byte-1.north west) -- (byte-1.south west)
    node [anchor=east,right=of byte-1,right,font=\scriptsize\ttfamily] {\textcolor{typeColor}{= 4}};
  \draw (byte-2.north west) -- (byte-2.south west)
    node [anchor=east,right=of byte-2,right,font=\scriptsize\ttfamily] {\textcolor{statusColor}{= 3}};
  \draw (byte-3.north west) -- (byte-3.south west)
    node [anchor=east,right=of byte-3,right,font=\scriptsize\ttfamily] {\textcolor{nameColor}{= 5}};
  \draw (byte-4.north west) -- (byte-4.south west)
    node [anchor=east,right=of byte-4,right,font=\scriptsize\ttfamily] {\textcolor{versionColor}{= 5}};
  \draw (byte-5.north west) -- (byte-5.south west)
    node [anchor=east,right=of byte-5,right,font=\scriptsize\ttfamily] {\textcolor{copyrightColor}{= 13}};
  \draw (byte-6.north west) -- (byte-6.south west)
    node [anchor=east,right=of byte-6,right,font=\scriptsize\ttfamily] {\textcolor{otherColor}{= 10}};
    
\end{tikzpicture}
\label{figure:Info2Seg2StrEx01SegLen}

\vspace{1em}
For a total of 65 bytes.

\end{minipage}


\section{\texttt{Info} Example 2}
\input{Info2StrEx02.tex}

\section{\texttt{Info} Example 3}
\begin{minipage}{\linewidth}
  
\begin{Verbatim}
Dwm::What::Info(\textcolor{typeColor}{"X"}, \textcolor{statusColor}{"rel"}, \textcolor{nameColor}{"myprog"}, \textcolor{versionColor}{"2.0.0"}, \textcolor{copyrightColor}{"ACME 2025"}, \textcolor{otherColor}{"Wile E. Coyote"})
\end{Verbatim}
\noindent\begin{tikzpicture}[start chain = A going right]
\node[draw=none,anchor=north west,yshift=1em,on chain=A] {\em compilation produces compile-time string};
\node[draw, single arrow,
  minimum height=10mm, minimum width=15mm,
  single arrow head extend=2mm,
  anchor=east, rotate=-90,xshift=-1em,on chain=A] {};
\end{tikzpicture}

\setcounter{lastbyte}{1}
\begin{tikzpicture}[
    start chain = A going right,
    node distance = 0pt, % Ensure no gap between bit nodes
    byte node/.style = {draw=lightgray, minimum height=20pt, minimum width=(\textwidth)/56, text width=\dimexpr(\textwidth)/56\relax, text height=16pt, text depth=4pt, font=\ttfamily, anchor=south, align=center, inner sep=0pt, outer sep=0pt, on chain=A}
]
  \def\mypfxstring{{@},{(},{\#},{)}}
  \def\mydelimstring{0xC2,0xA0}
  \def\mysegonestring{X}
  \def\mysegtwostring{r,e,l}
  \def\mysegthreestring{m,y,p,r,o,g}
  \def\mysegfourstring{2,.,0,.,0}
  \def\mysegfivestring{A,C,M,E,{ },2,0,2,5}
  \def\mysegsixstring{W,i,l,e,{ },E,.,{ },C,o,y,o,t,e}

  \foreach \byte [count=\i] in \mypfxstring {
    \node[byte node, fill=white] (byte-\arabic{lastbyte}) {\byte};
    \stepcounter{lastbyte}
  }
  \draw [decorate, decoration={brace, mirror, amplitude=10pt}]
    (byte-1.south west) -- (byte-4.south east)
    node [below=8pt, midway, font=\scriptsize] {\texttt{nth(0)}};
  \draw [decorate, decoration={brace, amplitude=10pt}]
    (byte-1.north west) -- (byte-4.north east)
    node [above=8pt, midway] {@(\#)};
  \foreach \byte [count=\i] in \mydelimstring {
    \node[byte node, text=delimColor, fill=white, font=\scriptsize\ttfamily, text depth=0pt] (byte-\arabic{lastbyte}) {\rotatebox{90}{\byte}};
    \stepcounter{lastbyte}
  }
  \draw [decorate, decoration={brace, amplitude=10pt}]
  (byte-5.north west) -- (byte-6.north east)
  node [above=8pt, midway] {\textvisiblespace};

  \foreach \byte [count=\i] in \mysegonestring {
    \node[byte node, fill=typeColor!10, text=typeColor, text depth=0pt] (byte-\arabic{lastbyte}) {\byte};
    \stepcounter{lastbyte}
  }
  \draw [decorate, decoration={brace, mirror, amplitude=10pt}]
    (byte-7.south west) -- (byte-7.south east)
    node [draw=none,below=8pt, midway, font=\scriptsize] {\texttt{nth(1)}}
    node [draw=none, below=18pt, midway, font=\scriptsize, text=typeColor] {\texttt{type()}};
  \draw [decorate, decoration={brace, amplitude=10pt}]
    (byte-7.north west) -- (byte-7.north east)
  node [above=8pt, midway] {X};
  \foreach \byte [count=\i] in \mydelimstring {
    \node[byte node, text=delimColor, fill=white, font=\scriptsize\ttfamily, text depth=0pt] (byte-\arabic{lastbyte}) {\rotatebox{90}{\byte}};
    \stepcounter{lastbyte}
  }
  \draw [decorate, decoration={brace, amplitude=10pt}]
  (byte-8.north west) -- (byte-9.north east)
  node [above=8pt, midway] {\textvisiblespace};

  \foreach \byte [count=\i] in \mysegtwostring {
    \node[byte node, fill=statusColor!10, text=statusColor, text depth=0pt] (byte-\arabic{lastbyte}) {\byte};
    \stepcounter{lastbyte}
  }
  \draw [decorate, decoration={brace, mirror, amplitude=10pt}]
    (byte-10.south west) -- (byte-12.south east)
    node [below=8pt, midway, font=\scriptsize] {\texttt{nth(2)}}
    node [below=18pt, midway, font=\scriptsize, text=statusColor] {\texttt{status()}};
  \draw [decorate, decoration={brace, amplitude=10pt}]
  (byte-10.north west) -- (byte-12.north east)
  node [above=8pt, midway] {rel};
  \foreach \byte [count=\i] in \mydelimstring {
    \node[byte node, text=delimColor, fill=white, font=\scriptsize\ttfamily, text depth=0pt] (byte-\arabic{lastbyte}) {\rotatebox{90}{\byte}};
    \stepcounter{lastbyte}    
  }
  \draw [decorate, decoration={brace, amplitude=10pt}]
  (byte-13.north west) -- (byte-14.north east)
  node [above=8pt, midway] {\textvisiblespace};

  \foreach \byte [count=\i] in \mysegthreestring {
    \node[byte node, fill=nameColor!10, text=nameColor] (byte-\arabic{lastbyte}) {\byte};
    \stepcounter{lastbyte}
  }
  \draw [decorate, decoration={brace, mirror, amplitude=10pt}]
    (byte-15.south west) -- (byte-20.south east)
    node [below=8pt, midway, font=\scriptsize] {\texttt{nth(3)}}
    node [below=18pt, midway, font=\scriptsize, text=nameColor] {\texttt{name()}};
  \draw [decorate, decoration={brace, amplitude=10pt}]
    (byte-15.north west) -- (byte-20.north east)
    node [above=8pt, midway] {myprog};
  \foreach \byte [count=\i] in \mydelimstring {
    \node[byte node, text=delimColor, fill=white, font=\scriptsize\ttfamily, text depth=0pt] (byte-\arabic{lastbyte}) {\rotatebox{90}{\byte}};
    \stepcounter{lastbyte}    
  }
  \draw [decorate, decoration={brace, amplitude=10pt}]
  (byte-21.north west) -- (byte-22.north east)
  node [above=8pt, midway] {\textvisiblespace};

  \foreach \byte [count=\i] in \mysegfourstring {
    \node[byte node, fill=versionColor!10, text=versionColor] (byte-\arabic{lastbyte}) {\byte};
    \stepcounter{lastbyte}
  }
  \draw [decorate, decoration={brace, mirror, amplitude=10pt}]
    (byte-23.south west) -- (byte-27.south east)
    node [below=8pt, midway, font=\scriptsize] {\texttt{nth(4)}}
    node [below=18pt, midway, font=\scriptsize, text=versionColor] {\texttt{version()}};
  \draw [decorate, decoration={brace, amplitude=10pt}]
    (byte-23.north west) -- (byte-27.north east)
    node [above=8pt, midway] {2.0.0};
  \foreach \byte [count=\i] in \mydelimstring {
    \node[byte node, text=delimColor, fill=white, font=\scriptsize\ttfamily, text depth=0pt] (byte-\arabic{lastbyte}) {\rotatebox{90}{\byte}};
    \stepcounter{lastbyte}    
  }
  \draw [decorate, decoration={brace, amplitude=10pt}]
  (byte-28.north west) -- (byte-29.north east)
  node [above=8pt, midway] {\textvisiblespace};

  \foreach \byte [count=\i] in \mysegfivestring {
    \node[byte node, fill=copyrightColor!10, text=copyrightColor] (byte-\arabic{lastbyte}) {\byte};
    \stepcounter{lastbyte}
  }
  \draw [decorate, decoration={brace, mirror, amplitude=10pt}]
    (byte-30.south west) -- (byte-38.south east)
    node [below=8pt, midway, font=\scriptsize] {\texttt{nth(5)}}
    node [below=18pt, midway, font=\scriptsize, text=copyrightColor] {\texttt{copyright()}};
  \draw [decorate, decoration={brace, amplitude=10pt}]
    (byte-30.north west) -- (byte-38.north east)
    node [above=8pt, midway] {ACME 2025};
  \foreach \byte [count=\i] in \mydelimstring {
    \node[byte node, text=delimColor, fill=white, font=\scriptsize\ttfamily, text depth=0pt] (byte-\arabic{lastbyte}) {\rotatebox{90}{\byte}};
    \stepcounter{lastbyte}    
  }
  \draw [decorate, decoration={brace, amplitude=10pt}]
  (byte-39.north west) -- (byte-40.north east)
  node [above=8pt, midway] {\textvisiblespace};

  \foreach \byte [count=\i] in \mysegsixstring {
    \node[byte node, fill=otherColor!10, text=otherColor] (byte-\arabic{lastbyte}) {\byte};
    \stepcounter{lastbyte}
  }
  \draw [decorate, decoration={brace, mirror, amplitude=10pt}]
    (byte-41.south west) -- (byte-54.south east)
    node [below=8pt, midway, font=\scriptsize] {\texttt{nth(6)}}
    node [below=18pt, midway, font=\scriptsize, text=otherColor] {\texttt{other()}};
  \draw [decorate, decoration={brace, amplitude=10pt}]
    (byte-41.north west) -- (byte-54.north east)
    node [above=8pt, midway] {Wile E. Coyote};

  \node[byte node, fill=lightgray!50] (termnull) {$\emptyset$};

  \draw [decorate, decoration={brace,mirror,amplitude=20pt,raise=24pt}]
    (byte-7.south west) -- (byte-54.south east)
    node [below=40pt, midway, font=\scriptsize,text=blue] {\texttt{data\_view()}};

  \draw [decorate, decoration={brace,mirror,amplitude=20pt,raise=48pt,aspect=0.45}]
    (byte-1.south west) -- (byte-54.south east)
    node [below=64pt, midway, xshift=-2.2em, font=\scriptsize] {\texttt{view()}};
\end{tikzpicture}
\label{figure:Info2Seg2StrEx03}

\begin{tikzpicture}[start chain = A going right]
\node[draw=none,minimum height=3em,anchor=north west,yshift=2em,on chain=A] {\em and an array of 7 bytes for segment lengths};
\end{tikzpicture}

\setcounter{lastbyte}{0}

\begin{tikzpicture}[
    start chain = A going below,
    node distance = 0pt, % Ensure no gap between bit nodes
    byte node/.style = {draw=lightgray, minimum height=(\textwidth)/54, minimum width=20pt, text width=16pt, text height=\dimexpr(\textwidth)/54\relax, text depth=2pt, font=\ttfamily, anchor=south, align=center, inner sep=0pt, outer sep=0pt, on chain=A}
]
    
  \def\mysegonestring{0x04,0x01,0x03,0x06,0x05,0x09,0x0E}

  \foreach \byte [count=\i from 0] in \mysegonestring {
    \node[byte node, font=\scriptsize\ttfamily, text depth=2pt] (byte-\arabic{lastbyte}) {\byte};
    \stepcounter{lastbyte}
  }
  \foreach \byte [count=\i from 0] in \mysegonestring {
    \draw [decorate, decoration={brace, mirror, amplitude=10pt}]
      (byte-\i.north west) -- (byte-\i.south west)
      node [anchor=west, midway, left, xshift=-1em,font=\scriptsize] (seglen-\i) {\texttt{nth(\i).size()}};
  }
  \draw (byte-1.north west) -- (byte-1.south west)
    node [anchor=west,left=of seglen-1,left,font=\scriptsize\ttfamily] {\textcolor{typeColor}{type().size(),}};
  \draw (byte-2.north west) -- (byte-2.south west)
    node [anchor=west,left=of seglen-2,left,font=\scriptsize\ttfamily] {\textcolor{statusColor}{status().size(),}};
  \draw (byte-3.north west) -- (byte-3.south west)
    node [anchor=west,left=of seglen-3,left,font=\scriptsize\ttfamily] {\textcolor{nameColor}{name().size(),}};
  \draw (byte-4.north west) -- (byte-4.south west)
    node [anchor=west,left=of seglen-4,left,font=\scriptsize\ttfamily] {\textcolor{versionColor}{version().size(),}};
  \draw (byte-5.north west) -- (byte-5.south west)
    node [anchor=west,left=of seglen-5,left,font=\scriptsize\ttfamily] {\textcolor{copyrightColor}{copyright().size(),}};
  \draw (byte-6.north west) -- (byte-6.south west)
    node [anchor=west,left=of seglen-6,left,font=\scriptsize\ttfamily] {\textcolor{otherColor}{other().size(),}};

  \draw (byte-0.north west) -- (byte-0.south west)
    node [anchor=east,right=of byte-0,right,font=\scriptsize\ttfamily] {\textcolor{black}{= 4}};
  \draw (byte-1.north west) -- (byte-1.south west)
    node [anchor=east,right=of byte-1,right,font=\scriptsize\ttfamily] {\textcolor{typeColor}{= 1}};
  \draw (byte-2.north west) -- (byte-2.south west)
    node [anchor=east,right=of byte-2,right,font=\scriptsize\ttfamily] {\textcolor{statusColor}{= 3}};
  \draw (byte-3.north west) -- (byte-3.south west)
    node [anchor=east,right=of byte-3,right,font=\scriptsize\ttfamily] {\textcolor{nameColor}{= 6}};
  \draw (byte-4.north west) -- (byte-4.south west)
    node [anchor=east,right=of byte-4,right,font=\scriptsize\ttfamily] {\textcolor{versionColor}{= 5}};
  \draw (byte-5.north west) -- (byte-5.south west)
    node [anchor=east,right=of byte-5,right,font=\scriptsize\ttfamily] {\textcolor{copyrightColor}{= 9}};
  \draw (byte-6.north west) -- (byte-6.south west)
    node [anchor=east,right=of byte-6,right,font=\scriptsize\ttfamily] {\textcolor{otherColor}{= 14}};

\end{tikzpicture}
\label{figure:Info2Seg2StrEx03SegLen}

\vspace{1em}
For a total of 62 bytes.

\end{minipage}


\section{\texttt{Info} Example 4}
If your toolchain and editor are UTF-8 aware, you can just paste
emoji directly into the string literal arguments for the
\texttt{Dwm::What::Info} constructor.

\begin{Verbatim}
Dwm::What::Info("\emoji{bomb}", "\emoji{collision}", \textcolor{nameColor}{"coyote"}, \textcolor{versionColor}{"9.4.2"}, \textcolor{copyrightColor}{"ACME 2025"}, \textcolor{otherColor}{"\emoji{feather} Beep Beep!"})
\end{Verbatim}
\noindent\begin{tikzpicture}[start chain = A going right]
\node[draw=none,anchor=north west,yshift=1em,on chain=A] {\em compilation produces compile-time string};
\node[draw, single arrow,
  minimum height=10mm, minimum width=15mm,
  single arrow head extend=2mm,
  anchor=east, rotate=-90,xshift=-1em,on chain=A] {};
\end{tikzpicture}

\setcounter{lastbyte}{1}
\begin{tikzpicture}[
    start chain = A going right,
    node distance = 0pt, % Ensure no gap between bit nodes
    byte node/.style = {draw=lightgray, minimum height=20pt, minimum width=(\textwidth)/56, text width=\dimexpr(\textwidth)/56\relax, text height=16pt, text depth=4pt, font=\ttfamily, anchor=south, align=center, inner sep=0pt, outer sep=0pt, on chain=A}
]
  \def\mypfxstring{{@},{(},{\#},{)}}
  \def\mydelimstring{0xC2,0xA0}
  \def\mysegonestring{0xF0,0x9F,0x92,0xA3}
  \def\mysegtwostring{0xF0,0x9F,0x92,0xA5}
  \def\mysegthreestring{c,o,y,o,t,e}
  \def\mysegfourstring{9,.,4,.,2}
  \def\mysegfivestring{A,C,M,E,{ },2,0,2,5}
  \def\mysegsixstringa{0xF0,0x9F,0xAA,0xB6}
  \def\mysegsixstringb{{ },B,e,e,p,{ },B,e,e,p,!}

  \foreach \byte [count=\i] in \mypfxstring {
    \node[byte node, fill=white] (byte-\arabic{lastbyte}) {\byte};
    \stepcounter{lastbyte}
  }
  \draw [decorate, decoration={brace, mirror, amplitude=10pt}]
    (byte-1.south west) -- (byte-4.south east)
    node [below=8pt, midway, font=\scriptsize] {\texttt{nth(0)}};
  \draw [decorate, decoration={brace, amplitude=10pt}]
    (byte-1.north west) -- (byte-4.north east)
    node [above=8pt, midway] {@(\#)};
  \foreach \byte [count=\i] in \mydelimstring {
    \node[byte node, text=delimColor, fill=white, font=\scriptsize\ttfamily, text depth=0pt] (byte-\arabic{lastbyte}) {\rotatebox{90}{\byte}};
    \stepcounter{lastbyte}
  }
  \draw [decorate, decoration={brace, amplitude=10pt}]
  (byte-5.north west) -- (byte-6.north east)
  node [above=8pt, midway] {\textvisiblespace};

  \foreach \byte [count=\i] in \mysegonestring {
    \node[byte node, fill=typeColor!10, text=typeColor, text depth=0pt,font=\scriptsize\ttfamily] (byte-\arabic{lastbyte}) {\rotatebox{90}{\byte}};
    \stepcounter{lastbyte}
  }
  \draw [decorate, decoration={brace, mirror, amplitude=10pt}]
    (byte-7.south west) -- (byte-10.south east)
    node [draw=none,below=8pt, midway, font=\scriptsize] {\texttt{nth(1)}}
    node [draw=none, below=18pt, midway, font=\scriptsize, text=typeColor] {\texttt{type()}};
  \draw [decorate, decoration={brace, amplitude=10pt}]
    (byte-7.north west) -- (byte-10.north east)
  node [above=8pt, midway] {\emoji{bomb}};
  \foreach \byte [count=\i] in \mydelimstring {
    \node[byte node, text=delimColor, fill=white, font=\scriptsize\ttfamily, text depth=0pt] (byte-\arabic{lastbyte}) {\rotatebox{90}{\byte}};
    \stepcounter{lastbyte}
  }
  \draw [decorate, decoration={brace, amplitude=10pt}]
  (byte-11.north west) -- (byte-12.north east)
  node [above=8pt, midway] {\textvisiblespace};

  \foreach \byte [count=\i] in \mysegtwostring {
    \node[byte node, fill=statusColor!10, ,font=\scriptsize\ttfamily, text=statusColor, text depth=0pt] (byte-\arabic{lastbyte}) {\rotatebox{90}{\byte}};
    \stepcounter{lastbyte}
  }
  \draw [decorate, decoration={brace, mirror, amplitude=10pt}]
    (byte-13.south west) -- (byte-16.south east)
    node [below=8pt, midway, font=\scriptsize] {\texttt{nth(2)}}
    node [below=18pt, midway, font=\scriptsize, text=statusColor] {\texttt{status()}};
  \draw [decorate, decoration={brace, amplitude=10pt}]
  (byte-13.north west) -- (byte-16.north east)
  node [above=8pt, midway] {\emoji{collision}};
  \foreach \byte [count=\i] in \mydelimstring {
    \node[byte node, text=delimColor, fill=white, font=\scriptsize\ttfamily, text depth=0pt] (byte-\arabic{lastbyte}) {\rotatebox{90}{\byte}};
    \stepcounter{lastbyte}    
  }
  \draw [decorate, decoration={brace, amplitude=10pt}]
  (byte-17.north west) -- (byte-18.north east)
  node [above=8pt, midway] {\textvisiblespace};

  \foreach \byte [count=\i] in \mysegthreestring {
    \node[byte node, fill=nameColor!10, text=nameColor] (byte-\arabic{lastbyte}) {\byte};
    \stepcounter{lastbyte}
  }
  \draw [decorate, decoration={brace, mirror, amplitude=10pt}]
    (byte-19.south west) -- (byte-24.south east)
    node [below=8pt, midway, font=\scriptsize] {\texttt{nth(3)}}
    node [below=18pt, midway, font=\scriptsize, text=nameColor] {\texttt{name()}};
  \draw [decorate, decoration={brace, amplitude=10pt}]
    (byte-19.north west) -- (byte-24.north east)
    node [above=8pt, midway] {coyote};
  \foreach \byte [count=\i] in \mydelimstring {
    \node[byte node, text=delimColor, fill=white, font=\scriptsize\ttfamily, text depth=0pt] (byte-\arabic{lastbyte}) {\rotatebox{90}{\byte}};
    \stepcounter{lastbyte}    
  }
  \draw [decorate, decoration={brace, amplitude=10pt}]
  (byte-25.north west) -- (byte-26.north east)
  node [above=8pt, midway] {\textvisiblespace};

  \foreach \byte [count=\i] in \mysegfourstring {
    \node[byte node, fill=versionColor!10, text=versionColor] (byte-\arabic{lastbyte}) {\byte};
    \stepcounter{lastbyte}
  }
  \draw [decorate, decoration={brace, mirror, amplitude=10pt}]
    (byte-27.south west) -- (byte-31.south east)
    node [below=8pt, midway, font=\scriptsize] {\texttt{nth(4)}}
    node [below=18pt, midway, font=\scriptsize, text=versionColor] {\texttt{version()}};
  \draw [decorate, decoration={brace, amplitude=10pt}]
    (byte-27.north west) -- (byte-31.north east)
    node [above=8pt, midway] {9.4.2};
  \foreach \byte [count=\i] in \mydelimstring {
    \node[byte node, text=delimColor, fill=white, font=\scriptsize\ttfamily, text depth=0pt] (byte-\arabic{lastbyte}) {\rotatebox{90}{\byte}};
    \stepcounter{lastbyte}    
  }
  \draw [decorate, decoration={brace, amplitude=10pt}]
  (byte-32.north west) -- (byte-33.north east)
  node [above=8pt, midway] {\textvisiblespace};

  \foreach \byte [count=\i] in \mysegfivestring {
    \node[byte node, fill=copyrightColor!10, text=copyrightColor] (byte-\arabic{lastbyte}) {\byte};
    \stepcounter{lastbyte}
  }
  \draw [decorate, decoration={brace, mirror, amplitude=10pt}]
    (byte-34.south west) -- (byte-42.south east)
    node [below=8pt, midway, font=\scriptsize] {\texttt{nth(5)}}
    node [below=18pt, midway, font=\scriptsize, text=copyrightColor] {\texttt{copyright()}};
  \draw [decorate, decoration={brace, amplitude=10pt}]
    (byte-34.north west) -- (byte-42.north east)
    node [above=8pt, midway] {ACME 2025};
  \foreach \byte [count=\i] in \mydelimstring {
    \node[byte node, text=delimColor, fill=white, font=\scriptsize\ttfamily, text depth=0pt] (byte-\arabic{lastbyte}) {\rotatebox{90}{\byte}};
    \stepcounter{lastbyte}    
  }
  \draw [decorate, decoration={brace, amplitude=10pt}]
  (byte-43.north west) -- (byte-44.north east)
  node [above=8pt, midway] {\textvisiblespace};

  \foreach \byte [count=\i] in \mysegsixstringa {
    \node[byte node, fill=otherColor!10, text=otherColor, font=\scriptsize\ttfamily, text depth=0pt] (byte-\arabic{lastbyte}) {\rotatebox{90}{\byte}};
    \stepcounter{lastbyte}
  }
  \foreach \byte [count=\i] in \mysegsixstringb {
    \node[byte node, fill=otherColor!10, text=otherColor] (byte-\arabic{lastbyte}) {\byte};
    \stepcounter{lastbyte}
  }
  \draw [decorate, decoration={brace, mirror, amplitude=10pt}]
    (byte-45.south west) -- (byte-59.south east)
    node [below=8pt, midway, font=\scriptsize] {\texttt{nth(6)}}
    node [below=18pt, midway, font=\scriptsize, text=otherColor] {\texttt{other()}};
  \draw [decorate, decoration={brace, amplitude=10pt}]
    (byte-45.north west) -- (byte-59.north east)
    node [above=8pt, midway] {\emoji{feather} Beep Beep!};

  \node[byte node, fill=lightgray!50] (termnull) {$\emptyset$};

  \draw [decorate, decoration={brace,mirror,amplitude=20pt,raise=24pt}]
    (byte-7.south west) -- (byte-59.south east)
    node [below=40pt, midway, font=\scriptsize,text=blue] {\texttt{data\_view()}};

  \draw [decorate, decoration={brace,mirror,amplitude=20pt,raise=48pt,aspect=0.45}]
    (byte-1.south west) -- (byte-59.south east)
    node [below=64pt, midway, xshift=-2.2em, font=\scriptsize] {\texttt{view()}};
\end{tikzpicture}
\label{figure:Info2Seg2StrEx04}

\begin{tikzpicture}[start chain = A going right]
\node[draw=none,minimum height=3em,anchor=north west,yshift=2em,on chain=A] {\em and an array of 7 bytes for segment lengths};
\end{tikzpicture}

\setcounter{lastbyte}{0}

\begin{tikzpicture}[
    start chain = A going below,
    node distance = 0pt, % Ensure no gap between bit nodes
    byte node/.style = {draw=lightgray, minimum height=(\textwidth)/54, minimum width=20pt, text width=16pt, text height=\dimexpr(\textwidth)/54\relax, text depth=2pt, font=\ttfamily, anchor=south, align=center, inner sep=0pt, outer sep=0pt, on chain=A}
]
    
  \def\mysegonestring{0x04,0x04,0x04,0x06,0x05,0x09,0x0F}

  \foreach \byte [count=\i from 0] in \mysegonestring {
    \node[byte node, font=\scriptsize\ttfamily, text depth=2pt] (byte-\arabic{lastbyte}) {\byte};
    \stepcounter{lastbyte}
  }
  \foreach \byte [count=\i from 0] in \mysegonestring {
    \draw [decorate, decoration={brace, mirror, amplitude=10pt}]
      (byte-\i.north west) -- (byte-\i.south west)
      node [anchor=west, midway, left, xshift=-1em,font=\scriptsize] (seglen-\i) {\texttt{nth(\i).size()}};
  }
  \draw (byte-1.north west) -- (byte-1.south west)
    node [anchor=west,left=of seglen-1,left,font=\scriptsize\ttfamily] {\textcolor{typeColor}{type().size(),}};
  \draw (byte-2.north west) -- (byte-2.south west)
    node [anchor=west,left=of seglen-2,left,font=\scriptsize\ttfamily] {\textcolor{statusColor}{status().size(),}};
  \draw (byte-3.north west) -- (byte-3.south west)
    node [anchor=west,left=of seglen-3,left,font=\scriptsize\ttfamily] {\textcolor{nameColor}{name().size(),}};
  \draw (byte-4.north west) -- (byte-4.south west)
    node [anchor=west,left=of seglen-4,left,font=\scriptsize\ttfamily] {\textcolor{versionColor}{version().size(),}};
  \draw (byte-5.north west) -- (byte-5.south west)
    node [anchor=west,left=of seglen-5,left,font=\scriptsize\ttfamily] {\textcolor{copyrightColor}{copyright().size(),}};
  \draw (byte-6.north west) -- (byte-6.south west)
    node [anchor=west,left=of seglen-6,left,font=\scriptsize\ttfamily] {\textcolor{otherColor}{other().size(),}};

  \draw (byte-0.north west) -- (byte-0.south west)
    node [anchor=east,right=of byte-0,right,font=\scriptsize\ttfamily] {\textcolor{black}{= 4}};
  \draw (byte-1.north west) -- (byte-1.south west)
    node [anchor=east,right=of byte-1,right,font=\scriptsize\ttfamily] {\textcolor{typeColor}{= 4}};
  \draw (byte-2.north west) -- (byte-2.south west)
    node [anchor=east,right=of byte-2,right,font=\scriptsize\ttfamily] {\textcolor{statusColor}{= 4}};
  \draw (byte-3.north west) -- (byte-3.south west)
    node [anchor=east,right=of byte-3,right,font=\scriptsize\ttfamily] {\textcolor{nameColor}{= 6}};
  \draw (byte-4.north west) -- (byte-4.south west)
    node [anchor=east,right=of byte-4,right,font=\scriptsize\ttfamily] {\textcolor{versionColor}{= 5}};
  \draw (byte-5.north west) -- (byte-5.south west)
    node [anchor=east,right=of byte-5,right,font=\scriptsize\ttfamily] {\textcolor{copyrightColor}{= 9}};
  \draw (byte-6.north west) -- (byte-6.south west)
    node [anchor=east,right=of byte-6,right,font=\scriptsize\ttfamily] {\textcolor{otherColor}{= 15}};

\end{tikzpicture}
\label{figure:Info2Seg2StrEx04SegLen}

\vspace{1em}
For a total of 67 bytes.




\chapter{More \texttt{SegmentedLiteral} examples}

\section{\texttt{SegmentedLiteral} Example 1}

\input{SegLit2StrEx01.tex}

\section{\texttt{SegmentedLiteral} Example 2}

\begin{minipage}{\linewidth}
  
\begin{Verbatim}
Dwm::What::SegmentedLiteral(" ", "@(#)", \textcolor{typeColor}{"The"}, \textcolor{statusColor}{"cow"}, \textcolor{nameColor}{"jumped"}, \textcolor{versionColor}{"over"}, \textcolor{copyrightColor}{"the"}, \textcolor{otherColor}{"moon"})
\end{Verbatim}
\noindent\begin{tikzpicture}[start chain = A going right]
\node[draw=none,anchor=north west,yshift=1em,on chain=A] {\em compilation produces compile-time string};
\node[draw, single arrow,
  minimum height=10mm, minimum width=15mm,
  single arrow head extend=2mm,
  anchor=east, rotate=-90,xshift=-1em,on chain=A] {};
\end{tikzpicture}

\setcounter{lastbyte}{1}
\begin{center}
\begin{tikzpicture}[
    start chain = A going right,
    node distance = 0pt, % Ensure no gap between bit nodes
    byte node/.style = {draw=lightgray, minimum height=20pt, minimum width=(\textwidth)/56, text width=\dimexpr(\textwidth)/56\relax, text height=16pt, text depth=4pt, font=\ttfamily, align=center, inner sep=0pt, outer sep=0pt, on chain=A}
]
% anchor=south, 
  \def\mypfxstring{{@},{(},{\#},{)}}
  \def\mydelimstring{ }
  \def\mysegonestring{T,h,e}
  \def\mysegtwostring{c,o,w}
  \def\mysegthreestring{j,u,m,p,e,d}
  \def\mysegfourstring{o,v,e,r}
  \def\mysegfivestring{t,h,e}
  \def\mysegsixstring{m,o,o,n}

  \foreach \byte [count=\i] in \mypfxstring {
    \node[byte node, fill=white] (byte-\arabic{lastbyte}) {\byte};
    \stepcounter{lastbyte}
  }
  \draw [decorate, decoration={brace, mirror, amplitude=10pt}]
    (byte-1.south west) -- (byte-4.south east)
    node [below=8pt, midway, font=\scriptsize] {\texttt{nth(0)}};
  \draw [decorate, decoration={brace, amplitude=10pt}]
    (byte-1.north west) -- (byte-4.north east)
    node [above=8pt, midway] {@(\#)};
  \foreach \byte [count=\i] in \mydelimstring {
    \node[byte node, text=delimColor, fill=white, font=\scriptsize\ttfamily, text depth=4pt] (byte-\arabic{lastbyte}) {\rotatebox{90}{\byte}};
    \stepcounter{lastbyte}
  }
  \draw [decorate, decoration={brace, amplitude=10pt}]
  (byte-5.north west) -- (byte-5.north east)
  node [above=8pt, midway] { };

  \foreach \byte [count=\i] in \mysegonestring {
    \node[byte node, fill=typeColor!10, text=typeColor, font=\ttfamily, text depth=4pt] (byte-\arabic{lastbyte}) {\byte};
    \stepcounter{lastbyte}
  }
  \draw [decorate, decoration={brace, mirror, amplitude=10pt}]
    (byte-6.south west) -- (byte-8.south east)
    node [below=8pt, midway, font=\scriptsize] {\texttt{nth(1)}};
  \draw [decorate, decoration={brace, amplitude=10pt}]
    (byte-6.north west) -- (byte-8.north east)
    node [above=8pt, midway] {The};
  \foreach \byte [count=\i] in \mydelimstring {
    \node[byte node, text=delimColor, fill=white, font=\scriptsize\ttfamily, text depth=0pt] (byte-\arabic{lastbyte}) {\rotatebox{90}{\byte}};
    \stepcounter{lastbyte}
  }
  \draw [decorate, decoration={brace, amplitude=10pt}]
  (byte-9.north west) -- (byte-9.north east)
  node [above=8pt, midway] {\textvisiblespace};

  \foreach \byte [count=\i] in \mysegtwostring {
    \node[byte node, fill=statusColor!10, font=\ttfamily, text=statusColor] (byte-\arabic{lastbyte}) {\byte};
    \stepcounter{lastbyte}
  }
  \draw [decorate, decoration={brace, mirror, amplitude=10pt}]
    (byte-10.south west) -- (byte-12.south east)
    node [below=8pt, midway, font=\scriptsize] {\texttt{nth(2)}};
  \draw [decorate, decoration={brace, amplitude=10pt}]
  (byte-10.north west) -- (byte-12.north east)
  node [above=8pt, midway] {cow};
  \foreach \byte [count=\i] in \mydelimstring {
    \node[byte node, text=delimColor, fill=white, font=\scriptsize\ttfamily, text depth=0pt] (byte-\arabic{lastbyte}) {\rotatebox{90}{\byte}};
    \stepcounter{lastbyte}    
  }
  \draw [decorate, decoration={brace, amplitude=10pt}]
  (byte-13.north west) -- (byte-13.north east)
  node [above=8pt, midway] {\textvisiblespace};

  \foreach \byte [count=\i] in \mysegthreestring {
    \node[byte node, fill=nameColor!10, text=nameColor] (byte-\arabic{lastbyte}) {\byte};
    \stepcounter{lastbyte}
  }
  \draw [decorate, decoration={brace, mirror, amplitude=10pt}]
    (byte-14.south west) -- (byte-19.south east)
    node [below=8pt, midway, font=\scriptsize] {\texttt{nth(3)}};
  \draw [decorate, decoration={brace, amplitude=10pt}]
    (byte-14.north west) -- (byte-19.north east)
    node [above=8pt, midway] {jumped};
  \foreach \byte [count=\i] in \mydelimstring {
    \node[byte node, text=delimColor, fill=white, font=\scriptsize\ttfamily, text depth=0pt] (byte-\arabic{lastbyte}) {\rotatebox{90}{\byte}};
    \stepcounter{lastbyte}    
  }
  \draw [decorate, decoration={brace, amplitude=10pt}]
  (byte-20.north west) -- (byte-20.north east)
  node [above=8pt, midway] {\textvisiblespace};

  \foreach \byte [count=\i] in \mysegfourstring {
    \node[byte node, fill=versionColor!10, text=versionColor] (byte-\arabic{lastbyte}) {\byte};
    \stepcounter{lastbyte}
  }
  \draw [decorate, decoration={brace, mirror, amplitude=10pt}]
    (byte-21.south west) -- (byte-24.south east)
    node [below=8pt, midway, font=\scriptsize] {\texttt{nth(4)}};
  \draw [decorate, decoration={brace, amplitude=10pt}]
    (byte-21.north west) -- (byte-24.north east)
    node [above=8pt, midway] {over};
  \foreach \byte [count=\i] in \mydelimstring {
    \node[byte node, text=delimColor, fill=white, font=\scriptsize\ttfamily, text depth=0pt] (byte-\arabic{lastbyte}) {\rotatebox{90}{\byte}};
    \stepcounter{lastbyte}    
  }
  \draw [decorate, decoration={brace, amplitude=10pt}]
  (byte-25.north west) -- (byte-25.north east)
  node [above=8pt, midway] {\textvisiblespace};

  \foreach \byte [count=\i] in \mysegfivestring {
    \node[byte node, fill=copyrightColor!10, text=copyrightColor] (byte-\arabic{lastbyte}) {\byte};
    \stepcounter{lastbyte}
  }
  \draw [decorate, decoration={brace, mirror, amplitude=10pt}]
    (byte-26.south west) -- (byte-28.south east)
    node [below=8pt, midway, font=\scriptsize] {\texttt{nth(5)}};
  \draw [decorate, decoration={brace, amplitude=10pt}]
    (byte-26.north west) -- (byte-28.north east)
    node [above=8pt, midway] {the};
  \foreach \byte [count=\i] in \mydelimstring {
    \node[byte node, text=delimColor, fill=white, font=\scriptsize\ttfamily, text depth=0pt] (byte-\arabic{lastbyte}) {\rotatebox{90}{\byte}};
    \stepcounter{lastbyte}    
  }
  \draw [decorate, decoration={brace, amplitude=10pt}]
  (byte-29.north west) -- (byte-29.north east)
  node [above=8pt, midway] {\textvisiblespace};

  \foreach \byte [count=\i] in \mysegsixstring {
    \node[byte node, fill=otherColor!10, text=otherColor] (byte-\arabic{lastbyte}) {\byte};
    \stepcounter{lastbyte}
  }
  \draw [decorate, decoration={brace, mirror, amplitude=10pt}]
    (byte-30.south west) -- (byte-33.south east)
    node [below=8pt, midway, font=\scriptsize] {\texttt{nth(6)}};
  \draw [decorate, decoration={brace, amplitude=10pt}]
    (byte-30.north west) -- (byte-33.north east)
    node [above=8pt, midway] {moon};

  \node[byte node, fill=lightgray!50] (termnull) {$\emptyset$};

  \draw [decorate, decoration={brace,mirror,amplitude=20pt,raise=20pt}]
    (byte-1.south west) -- (byte-33.south east)
    node [below=40pt, midway, font=\scriptsize] {\texttt{view()}};
    
\end{tikzpicture}
\label{figure:SegLit2StrEx02}
\end{center}

\begin{tikzpicture}[start chain = A going right]
\node[draw=none,minimum height=3em,anchor=north west,yshift=2em,on chain=A] {\em and an array of 7 bytes for segment lengths};
\end{tikzpicture}

\setcounter{lastbyte}{0}
\begin{tikzpicture}[
    start chain = A going below,
    node distance = 0pt, % Ensure no gap between bit nodes
    byte node/.style = {draw=lightgray, minimum height=(\textwidth)/56, minimum width=20pt, text width=16pt, text height=\dimexpr(\textwidth)/56\relax, text depth=2pt, font=\ttfamily, anchor=south, align=center, inner sep=0pt, outer sep=0pt, on chain=A}
]
    
  \def\mysegonestring{0x04,0x03,0x03,0x06,0x04,0x03,0x04}

  \foreach \byte [count=\i from 0] in \mysegonestring {
    \node[byte node, font=\scriptsize\ttfamily, text depth=2pt] (byte-\arabic{lastbyte}) {\byte};
    \stepcounter{lastbyte}
  }
  \foreach \byte [count=\i from 0] in \mysegonestring {
    \draw [decorate, decoration={brace, mirror, amplitude=10pt}]
      (byte-\i.north west) -- (byte-\i.south west)
      node [anchor=west, midway, left, xshift=-1em,font=\scriptsize] (seglen-\i) {\texttt{nth(\i).size()}};
  }

  \draw (byte-0.north west) -- (byte-0.south west)
    node [anchor=east,right=of byte-0,right,font=\scriptsize\ttfamily] {\textcolor{black}{= 4}};
  \draw (byte-1.north west) -- (byte-1.south west)
    node [anchor=east,right=of byte-1,right,font=\scriptsize\ttfamily] {\textcolor{typeColor}{= 3}};
  \draw (byte-2.north west) -- (byte-2.south west)
    node [anchor=east,right=of byte-2,right,font=\scriptsize\ttfamily] {\textcolor{statusColor}{= 3}};
  \draw (byte-3.north west) -- (byte-3.south west)
    node [anchor=east,right=of byte-3,right,font=\scriptsize\ttfamily] {\textcolor{nameColor}{= 6}};
  \draw (byte-4.north west) -- (byte-4.south west)
    node [anchor=east,right=of byte-4,right,font=\scriptsize\ttfamily] {\textcolor{versionColor}{= 4}};
  \draw (byte-5.north west) -- (byte-5.south west)
    node [anchor=east,right=of byte-5,right,font=\scriptsize\ttfamily] {\textcolor{copyrightColor}{= 3}};
  \draw (byte-6.north west) -- (byte-6.south west)
    node [anchor=east,right=of byte-6,right,font=\scriptsize\ttfamily] {\textcolor{otherColor}{= 4}};
    
\end{tikzpicture}
\label{figure:SegLit2StrEx02SegLen}

\vspace{1em}
For a total of 41 bytes.

\end{minipage}


\section{\texttt{SegmentedLiteral} Example 3}
\begin{minipage}{\linewidth}

An empty delimiter is permitted.
\vspace{1em}

\input{SegLit2StrEx03.tex}
\end{minipage}


\chapter{Macros in \texttt{DwmWhatInfo.hh}}
There are many macros in \texttt{DwmWhatInfo.hh} that expand to UTF-8 string literals,
just so I can get emoji output.  For my own use, the following macros are commonly
used.

\begin{tabular}{llll}
  \textbf{macro} & \textbf{expands to (UTF-8)} & \textbf{Unicode Code Point} & \textbf{display} \\ \hline
  \texttt{DWM\_WHAT\_TYPE\_EXE}   & \texttt{"\textbackslash{xF0}\textbackslash{x9F}\textbackslash{xA4}\textbackslash{x96}"} & \texttt{U+1F916}   & \texttt{\emoji{robot}} \\
  \texttt{DWM\_WHAT\_TYPE\_LIB}   & \texttt{"\textbackslash{xF0}\textbackslash{x9F}\textbackslash{x93}\textbackslash{x9A}"} & \texttt{U+1F4DA}   & \texttt{\emoji{books}} \\
  \texttt{DWM\_WHAT\_TYPE\_HDR}   & \texttt{"\textbackslash{xEF}\textbackslash{xBC}\textbackslash{x83}"} & \texttt{U+FF03}    & \texttt{\#} \\
  \texttt{DWM\_WHAT\_TYPE\_DOC}   & \texttt{"\textbackslash{xF0}\textbackslash{x9F}\textbackslash{x93}\textbackslash{x84}"} & \texttt{U+1F4C4}   & \texttt{\emoji{page-facing-up}} \\
  \texttt{DWM\_WHAT\_STATUS\_REL} & \texttt{"\textbackslash{xE2}\textbackslash{x9C}\textbackslash{x85}"} & \texttt{U+2705}    & \texttt{\emoji{white-check-mark}} \\
  \texttt{DWM\_WHAT\_STATUS\_RC}  & \texttt{"\textbackslash{xF0}\textbackslash{x9F}\textbackslash{x91}\textbackslash{xB7}"} & \texttt{U+1F477}   & \texttt{\emoji{construction-worker}} \\
  \texttt{DWM\_WHAT\_STATUS\_DEV} & \texttt{"\textbackslash{xE2}\textbackslash{x9D}\textbackslash{x97}"} & \texttt{U+2757}    & \texttt{\emoji{heavy-exclamation-mark}} \\
  \texttt{DWM\_WHAT\_SYM\_COPYRIGHT} & \texttt{"\textbackslash{xC2}\textbackslash{xA9}"} & \texttt{U+00A9} & \texttt{\copyright} \\
\end{tabular}
  
  


% \emoji{robot}
% \emoji{ghost}
% \emoji{heavy-exclamation-mark}
% \emoji{books}
% \emoji{white-check-mark}
% \emoji{construction-worker}
% \emoji{page-facing-up}
% \emoji{copyright}

\end{document}
